\documentclass[11pt]{article}
\usepackage[a4paper,margin=21mm]{geometry}
\usepackage[no-math]{fontspec}
\setmainfont{Latin Modern Roman}
\newfontfamily\CJKfont{Noto Serif CJK SC}
\XeTeXlinebreaklocale "zh"
\XeTeXlinebreakskip=0pt plus 1pt
\usepackage{amsmath,amssymb,amsthm,amscd,graphicx,color}
\usepackage{xurl}
\usepackage[unicode,hidelinks]{hyperref}
\allowdisplaybreaks
\setlength\emergencystretch{3em}
\numberwithin{equation}{section}

\numberwithin{equation}{section}

\usepackage{amstext,amscd,array}

\usepackage{mathrsfs}
\usepackage{bm}
\usepackage[arrow,matrix,curve]{xy}

\DeclareMathOperator{\AdS}{AdS}
\DeclareMathOperator{\Sphere}{\mathbb{S}}
\DeclareMathOperator{\NW}{NW}
\DeclareMathOperator{\CW}{CW}
\DeclareMathOperator{\ad}{ad}
\let\S\Sphere
\DeclareMathOperator{\weyl}{{\sf W}} %% \weyl =
%% Weyl operator

\newcommand{\comp}{\boxplus} %% \comp =
%% composition
\newcommand{\compa}{\mkern+4mu\square\mkern-17.4mu{\mbox{
 \protect\raisebox{0.2ex}{$\ast$}}\mkern+5mu}} %% \compa =
%% composition style a
\newcommand{\compb}{\mkern+4mu\square\mkern-17.4mu{\mbox{
 \protect\raisebox{0.2ex}{$\bullet$}}\mkern+5mu}} %% \compb =
%% composition style b
\newcommand{\compc}{\mkern+4mu\square\mkern-17.4mu{\mbox{
 \protect\raisebox{0.2ex}{$\star$}}\mkern+5mu}} %% \compc =
%% composition style c
\newcommand{\cb}[2]{\bigl[#1\,,\,#2\bigr]} %% \cb{a}{b} =
%% [ a, b ]
\newcommand{\scb}[2]{\bigl[#1\,,\,#2\bigr]_\star} %% \scb{a}{b}
%% = [ a, b ]_star
\newcommand{\NO}{\mbox{$\substack{\circ\\\circ}$}} %% Normal
%% ordering
\newcommand{\NOa}{\mbox{$\substack{\ast\\\ast}$}} %% Normal
%% ordering alt A
\newcommand{\NOb}{\mbox{$\substack{\bullet\\\bullet}$}} %% Normal
%% ordering alt B

\def\nn{\nonumber}
\def\d{\partial}
\def\bfd{{\mbf\partial}}
\def\bfz{{\mbf z}}
\def\od{\overline{\partial}}
\def\oz{\overline{z}}
\def\obfd{\overline{\mbf\partial}}
\def\obfz{\overline{\mbf z}}
\def\tri{\triangleright}
\def\rbtri{\blacktriangleright}
\def\lbtri{\blacktriangleleft}
%\newcommand{\1}{\mathbb{1}}

\newcommand{\appx}{\approx}
\newcommand{\cD}{{\mathcal D}}

\newcommand{\sgn}{{\rm sgn}}
\newcommand{\eps}{\varepsilon}
\newcommand{\half}{\mbox{$\frac{1}{2}$}}
\newcommand{\behalf}{\frac{\beta}{2}}
\def\Dirac{{D\!\!\!\!/\,}} % Dirac operator

\newcommand{\C}{\complex}
\newcommand{\Z}{\zed}
\newcommand{\N}{\nat}
\newcommand{\F}{\mathbb{F}}

\def\one{\mathbb{1}}

\newcommand{\QED}{\hfill QED}
\newcommand{\mbf}[1]{{\boldsymbol {#1} }}
\def\ii{{{\rm i}}}
\def\dd{{\rm d}}
\def\DD{{\rm D}}
\def\CC{{\rm C}}
\def\Id{{\rm id}}
\def\P{{\sf P}}
\def\B{{\sf B}}
\def\a{{\sf a}}
\def\b{{\sf b}}
\def\sfd{{\sf d}}
\def\sfs{{\sf s}}
\def\sft{{\sf t}}
\def\sfi{{\sf i}}
\def\A{{\sf A}}
\def\V{{\sf V}}
\def\W{{\sf W}}
\def\K{{\sf K}}
\def\T{{\sf T}}
\def\X{{\sf X}}
\def\Y{{\sf Y}}
\def\Q{{\sf Q}}
\def\C{{\sf C}}
\def\J{{\sf J}}
\def\Diff{{\sf Diff}}
\def\ttp{{\tt p}}
\def\ttx{{\tt x}}
\def\CE{{\tt CE}}
\def\ellfrac{{\tfrac{m\,\pi}2\,x}}

\newcommand{\unit}{\mathbbm{1}} 			% identity map/matrix

\def\ma{{\mbf a}}
\def\mz{{\mbf z}}
\def\mw{{\mbf w}}
\def\mx{{\mbf x}}
\def\my{{\mbf y}}
\def\mk{{\mbf k}}
\def\mq{{\mbf q}}
\def\mbp{{\mbf p}}
\def\mD{{\mbf D}}
\def\mG{{\mbf G}}
\def\mdell{{\mbf\partial}}

\def\mff{{\mathfrak f}}
\def\mfn{{\mathfrak n}}
\def\mfg{{\mathfrak g}}
\def\mfh{{\mathfrak h}}
\def\mfs{{\mathfrak s}}
\newcommand{\fru}{\mathfrak{u}}				% frak-letters
\def\mcN{{\mathcal N}}
\def\mcS{{\mathcal S}}
\def\mbbV{{\mathbb V}}
\def\mcH{{\mathcal H}}
\def\mcF{{\mathcal F}}
\def\mcR{{\mathcal R}}

\def\Lie{{\mathcal L}}
\newcommand{\CV}{\mathcal{V}}
\newcommand{\CCV}{\mathscr{V}}
\newcommand{\CCC}{\mathscr{C}}
\newcommand{\CCG}{\mathscr{G}}
\newcommand{\CCB}{\mathscr{B}}
\newcommand{\CCA}{\mathscr{A}}
\newcommand{\CCW}{\mathscr{W}}
\newcommand{\CCP}{\mathscr{P}}
\newcommand{\CCF}{\mathscr{F}}
\newcommand{\CCM}{\mathscr{M}}
\newcommand{\CCK}{\mathscr{K}}
\newcommand{\CP}{\mathcal{P}}
\newcommand{\CX}{\mathcal{X}}
\newcommand{\CY}{\mathcal{Y}}
\newcommand{\CZ}{\mathcal{Z}}
\newcommand{\CK}{\mathcal{K}}
\newcommand{\CI}{\mathcal{I}}
\newcommand{\Dcal}{\mathcal{D}}
\newcommand{\Ccal}{\mathcal{C}}

\newcommand{\sfG}{\mathsf{G}}
\newcommand{\sfH}{\mathsf{H}}
\newcommand{\sfN}{\mathsf{N}}
\newcommand{\sfM}{\mathsf{M}}
\newcommand{\sfK}{\mathsf{K}}
\newcommand{\sfU}{\mathsf{U}}
\newcommand{\sfGL}{\mathsf{GL}}
\newcommand{\sfSL}{\mathsf{SL}}
\newcommand{\sfSU}{\mathsf{SU}}

\newcommand{\ttM}{\mathtt{M}}

\newcommand{\U}{{\rm U}}
\newcommand{\rmH}{{\rm H}}
\newcommand{\Pim}{{\mit\Pi}}
\newcommand{\End}{{\rm End}}
\newcommand{\Aut}{\operatorname{Aut}}
\newcommand{\Leb}{{\rm L}}
\newcommand{\Ind}{{\rm Ind}}
\newcommand{\Heis}{{\sf Heis}}
\newcommand{\ISO}{{\sf ISO}}
\newcommand{\SU}{{\sf SU}}
\newcommand{\vect}{{\sf Vect}}
\newcommand{\YM}{{\mathcal Y\!M}}
\newcommand{\tet}{\theta}
\newcommand{\cL}{{\mathcal L}}
\newcommand{\sine}{{\sf K}}
\newcommand{\rmB}{{\rm B}}
\newcommand{\lt}{{\rm lt}}
\newcommand{\rt}{{\rm rt}}

\newcommand{\mcoprod}{\mbox{$\coprod\limits_{x\in\R/\Z}$}}

\newcommand{\R}{\real}

\newcommand{\complex}{{\mathbb C}} %% complex numbers
\newcommand{\zed}{{\mathbb Z}} %% integers
\newcommand{\nat}{{\mathbb N}} %% naturals
\newcommand{\real}{{\mathbb R}} %% real numbers
\newcommand{\eucl}{{\mathbb E}}
\newcommand{\rat}{{\mathbb Q}} %% rational numbers
\newcommand{\mat}{{\mathbb M}} %% matrix algebra
\newcommand{\torus}{{\mathbb T}} %% torus
\newcommand{\quat}{{\mathbb H}} %% quaternions
\newcommand{\sphere}{{\mathbb S}} %% sphere
%\newcommand{\NO}{\,\mbox{$\circ\atop\circ$}\,} % Normal ordering
\newcommand{\id}{{1\!\!1}} %% identity matrix
\def\Dirac{{D\!\!\!\!/\,}} % Dirac operator
\def\semiprod{{\,\rhd\!\!\!<\,}} % crossed product
\def\alg{{\mathcal A}}
\def\balg{{\mathcal B}}
\def\hil{{\mathcal H}}
\def\Lcal{{\mathcal L}}
\def\Mcal{{\mathcal M}}
\def\Tcal{{\mathcal T}}
\def\ot{{\, \otimes\, }}
\def\ota{\otimes_{ A}}
\def\otc{\otimes_{\complexs}}
\def\otr{\otimes_{reals}}
\newif\ifold \oldtrue \def\new{\oldfalse}

\font\mathsm=cmmi9

\def\nn{\nonumber}
\newcommand{\tr}[1]{{\rm tr}\,#1}
\newcommand{\Tr}[1]{{\rm Tr}\,#1}
\newcommand{\sdet}[1]{\:{\rm SDet}\,#1}
\def\e{{\,\rm e}\,}
\newcommand{\rf}[1]{(\ref{#1})}
\newcommand{\non}{\nonumber \\*}
\newcommand{\vgr}{{\vec \nabla}}
\newcommand{\vx}{{\vec x}}

\def\const{{\rm const}}
\def\s{\sigma}
\def\vcl{\varphi_{\rm cl}}
\def\la{\left\langle}
\def\ra{\right\rancle}
\def\d{\partial}
\def\se{S_{\rm eff}}
\def\E{F}

\def\Ecal{E}

\newcommand{\fr}[2]{{\textstyle {#1 \over #2}}}

\newcommand{\fintf}[1]{\int \frac{d^{2d} #1}{(2\pi)^{2d}}}
\newcommand{\fintt}[1]{\int \frac{d^2 #1}{(2\pi)^2}}

\newcommand{\ph}{\phantom{ }}

\newcommand{\MM}{\bf \mathcal M}
\newcommand{\g}{\gamma}
\newcommand{\G}{\Gamma}
\newcommand{\rh}{\rho}
\newcommand{\sg}{\sigma}
\newcommand{\Sg}{\Sigma}
\newcommand{\pr}{\partial}
\newcommand{\tSg}{\tilde{\Sigma}}
\newcommand{\hr}{\hat{r}}
\newcommand{\hq}{\hat{q}}
\newcommand{\hk}{\hat{K}}
\newcommand{\ho}{\hat{\omega}}
\newcommand{\del}{\partial}
\newcommand{\ka}{ e}
\newcommand{\co}{{\mathcal O}}
\newcommand{\z}{\zeta}
\newcommand{\upa}{\uparrow}
\newcommand{\doa}{\downarrow}
\newcommand{\ZZ}{{\mathcal Z}}
\newcommand{\proj}{{\rm pr}}

\makeatletter
\newdimen\normalarrayskip % skip between lines
\newdimen\minarrayskip % minimal skip between lines
\normalarrayskip\baselineskip
\minarrayskip\jot
\newif\ifold \oldtrue \def\new{\oldfalse}
%
\def\arraymode{\ifold\relax\else\displaystyle\fi} % mode of array entries
\def\eqnumphantom{\phantom{(\theequation)}} % right phantom in eqnarray
\def\@arrayskip{\ifold\baselineskip\z@\lineskip\z@
 \else
 \baselineskip\minarrayskip\lineskip2\minarrayskip\fi}
%
\def\@arrayclassz{\ifcase \@lastchclass \@acolampacol \or
\@ampacol \or \or \or \@addamp \or
 \@acolampacol \or \@firstampfalse \@acol \fi
\edef\@preamble{\@preamble
 \ifcase \@chnum
 \hfil$\relax\arraymode\@sharp$\hfil
 \or $\relax\arraymode\@sharp$\hfil
 \or \hfil$\relax\arraymode\@sharp$\fi}}
%
\def\@array[#1]#2{\setbox\@arstrutbox=\hbox{\vrule
 height\arraystretch \ht\strutbox
 depth\arraystretch \dp\strutbox
 width\z@}\@mkpream{#2}\edef\@preamble{\halign \noexpand\@halignto
\bgroup \tabskip\z@ \@arstrut \@preamble \tabskip\z@ \cr}%
\let\@startpbox\@@startpbox \let\@endpbox\@@endpbox
 \if #1t\vtop \else \if#1b\vbox \else \vcenter \fi\fi
 \bgroup \let\par\relax
 \let\@sharp##\let\protect\relax
 \@arrayskip\@preamble}
\makeatother

\newcommand{\ov}{\overline}
\newcommand{\wt}{\widetilde}
\newcommand{\wg}{\wedge}
\newcommand{\vphi}{\varphi}
\newcommand{\p}{\partial}
\newcommand{\mc}{\mathcal}
\newcommand{\ms}{\mathscr}
\newcommand{\om}{\omega}
\newcommand{\de}{\delta}
\newcommand{\mn}{{\mu\nu}}
\newcommand{\dg}{\dagger}
\newcommand{\trm}{\textrm}
\newcommand{\tbf}{\textbf}
\newcommand{\wh}{\widehat}
\newcommand{\nab}{\nabla}
\newcommand{\fs}{\footnotesize}
\newcommand{\scr}{\scriptsize}

					% tilde letters and symbols
\def\tmfg{{\mfg}_e}
\def\tmfh{{\bar \mfh}}
\def\tH{{\bar H}}
\def\tF{{\bar F}}
\def\tP{{\tilde P}}
\def\tp{{\bar p}}
\def\tM{{\bar M}}
\def\tmcH{{\bar \mcH}}
\def\tmcF{{\bar \mcF}}
\def\As{{\APLstar}}
\def\tAs{{\,\bar \As \,}}
\def\ts{{\, \bar \star \,}}
\def\tmcR{{\bar \mcR}}
\def\moyal{{\, \mbox{{\tiny \FourStar}}\,}}

\def\uf{{\underline f}}
\def\ug{{\underline g}}
\def\uh{{\underline h}}
\def\uk{{\underline k}}
\def\uu{{\underline u}}
\def\uv{{\underline v}}
\def\uz{{\underline z}}

\def\FF{{\cal F}}
\def\al{\alpha}
\def\be{\beta}

\def \rv{\vec{r}}

\newcommand{\intps}{\int_{{\cal M}} \, \dd^{2d}x \ }
\def\cirp{\mathop{\bar\circ}}
\def\pphim{{\bar\phi\bar\phi}}

\newtheorem{Theorem}{Theorem}[section]
\newtheorem{Corollary}[Theorem]{Corollary}
\newtheorem{Lemma}[Theorem]{Lemma}
\newtheorem{Proposition}[Theorem]{Proposition}
 { \theoremstyle{definition}
\newtheorem{Definition}[Theorem]{Definition}
\newtheorem{Note}[Theorem]{Note}
\newtheorem{Example}[Theorem]{Example}
\newtheorem{Remark}[Theorem]{Remark} }

\newfontfamily\nativebbone{DejaVu Math TeX Gyre}
\ExplSyntaxOn
\NewDocumentCommand{\mathbbm}{m}{\str_if_eq:nnTF{#1}{1}{\mathord{\text{\upshape\nativebbone\char"1D7D9}}}{\PackageError{nativecompat}{Unsupported~mathbbm~argument}{Only~verified~digit~one~is~accepted}}}
\ExplSyntaxOff
\makeatletter\newlabel{thm:Takaiduality}{{2.9}{}{Original source locator}{}{}}\makeatother
\makeatletter\newlabel{sec:toruscrossedprod}{{3.4}{}{Original source locator}{}{}}\makeatother
\begin{document}
\begin{center}\Large Rank-one R\_z T-duals of compact three-dimensional almost abelian solvmanifolds are Morita-equivalent noncommutative two-torus bundles over the Mostow circle\end{center}
\noindent\textbf{Stable ID:} 1340\quad\textbf{Author:} Paolo Aschieri; Richard J. Szabo\par
\noindent\textbf{Work:} Topological T-Duality for Twisted Tori\par
\noindent\textbf{Source:} \url{https://sigma-journal.com/2021/012/}\par
\noindent\textbf{Version:} Official SIGMA 17 (2021) article 012, DOI 10.3842/SIGMA.2021.012; journal PDF and native TeX acquired 2026-09-30, exact bytes pinned by SHA256\par
\noindent\textbf{Rights:} CC BY 4.0. Authors retain copyright. \url{https://creativecommons.org/licenses/by/4.0/}. Reuse and typesetting adaptation; original language and mathematical content preserved.\par
\noindent\textbf{Difficulty (prior selection preserved):} {\CJKfont H2: The proof explicitly parameterizes the right-action cosets and reconstructs the induced lattice translations before decomposing the crossed product. It constructs an equivariant imprimitivity module to identify the nonsingular fibers with noncommutative two-tori and treats the exceptional zero-denominator fibers separately. The monodromy is then transported to a fractional-linear action and a Morita equivalence, supplying the geometric fiber description rather than only invoking abstract T-duality.}\par
\medskip\noindent\textbf{Editorial boundary note (not author text):} Selected central Theorem4.17 in d=3 for the nontrivial lattice-rational torsion-direction R\_z action of4.6, with R\_z intersect Lambda a lattice, complete period/monodromy/Mostow/Smith conventions and exceptional-fiber treatment. Full exact coset factorization, induced lattice action, crossed-product decomposition, explicit Green/equivariant Morita modules and noncommutative-torus identification are retained, as are the author monodromy formulas/calculation. General printed d-dimensional context remains literal and uncounted. Standard prior Green/Takai/Connes–Thom/imprimitivity/background stay original precise references; no new C*-machinery or editor proof. The only mathbbm expansion is the literal author unit macro with argument1, resolved by the approved finite double-struck-one named DejaVu Math TeX Gyre font (U+1D7D9), rejecting every other argument; no ordinary-one substitute. pdfsync debug and unavailable bbm import omitted. No formula/index/lattice-label correction. Unused bbding/wasysym imports also omitted after zero body calls to their star glyphs or dependent author macros. The indicator alone uses its explicitly upright named font even inside italic theorem text, avoiding missing italic-font substitution.\par
\noindent\textbf{External original reference thm:Takaiduality:} Prior Takai duality as originally stated in Theorem2.9, native1263–1271, with precise source reference Williams2007; not a missing new selected construction.\par
\noindent\textbf{External original reference sec:toruscrossedprod:} Original unselected torus/orbifold T-duality examples in3.4–3.5, native1902–2090; the selected Almost Abelian computations are supplied in full.\par
\medskip\hrule\medskip\noindent\textbf{Author text begins below.}\par
\setcounter{section}{1}
% Original source lines 621--754
\section{Crossed products and duality\label{sec:NAdef}}

In this section we summarise some of the mathematical tools that we
will use in this paper. A~good reference for the material covered in
the following is the book~\cite{Williams2007}. Throughout this paper,
all topological spaces are assumed to be second countable (hence
separable), locally compact and Hausdorff.

\subsection{Dynamical systems and their crossed products}

Let $\sfG$ be a locally compact group and let $X$ be a $\sfG$-space, i.e., a
topological space which is acted
upon
homeomorphically by $\sfG$; we denote the $\sfG$-action $\sfG\times
X\to X$ by $(\gamma,x)\mapsto\gamma\cdot x$. The pair $(X,\sfG)$ is
called a \emph{transformation group}.
 A related concept is that of a
\emph{dynamical system}, which is a triple $(\alg,\sfG,\alpha)$
consisting of an algebra $\alg$ and a locally compact group $\sfG$ acting
on~$\alg$ via a group homomorphism $\alpha\colon \sfG\to\Aut (\alg)$,
denoted $\gamma\mapsto(\alpha_\gamma\colon \alg\to\alg)$ for $\gamma\in
\sfG$. In topological
T-duality one usually requires $\alg$ to be a $C^*$-algebra, in which
case $(\alg,\sfG,\alpha)$ is called a \emph{$C^*$-dynamical
 system}. Two dynamical systems $(\alg,\sfG,\alpha)$ and $(\balg,\sfG,\beta)$ are
\emph{equivalent} if there is an algebra isomorphism
$\varphi\colon \alg\to\balg$ which intertwines the $\sfG$-actions: $\varphi\circ\alpha_\gamma =
\beta_\gamma\circ\varphi$ for all $\gamma\in \sfG$.

If $\alg$ is a commutative $C^*$-algebra, then we call $(\alg,\sfG,\alpha)$ a
\emph{commutative} dynamical system. In that case, by Gelfand duality
$\alg=C_0(X)$ is the algebra of continuous functions vanishing at
infinity on a
topological space $X$
equipped with the $\sfG$-action \smash{$\alpha_\gamma^\dag\big|_X$} for $\gamma\in \sfG$ under
the identification of points $x\in X$ with irreducible representations
of $C_0(X)$, which are one-dimensional and given by the point evaluation
maps ${\rm ev}_x\colon \alg\to\alg$ with ${\rm ev}_x(\mathtt{f})=\mathtt{f}(x)$; then
$(X,\sfG)$ is a~transformation group. Conversely, given a transformation
group $(X,\sfG)$, there is an associated commutative dynamical system
$(C_0(X), \sfG,\alpha)$, where $\alpha_\gamma(\mathtt{f})(x)=\mathtt{f}\big(\gamma^{-1}\cdot x\big)$ for $\gamma\in
\sfG$, $\mathtt{f}\in C_0(X)$ and $x\in X$. In other words, there is a one-to-one
correspondence between transformation groups and commutative $C^*$-dynamical
systems.

If the $C^*$-algebra $\alg$ is not commutative, then we call
$(\alg,\sfG,\alpha)$ a \emph{noncommutative} $C^*$-dynamical system.

As usual, it is more useful to work with a representation rather than the
abstract dynamical system itself. A \emph{covariant representation} of
a dynamical system $(\alg,\sfG,\alpha)$ in a $C^*$-algebra $\balg$ with
multiplier algebra ${\rm M}(\balg)$ is a
pair $(\Pim,U)$ consisting of a~homomorphism $\Pim\colon \alg\to {\rm M}(\balg)$ and a~unitary representation $U\colon \sfG\to{\rm M}(\balg)$, $\gamma\mapsto U_\gamma$ which satisfy the compatibility condition
\[
\Pim\big(\alpha_\gamma(a)\big) = U_\gamma \Pim(a) U_\gamma^{-1} ,
\]
for all $\gamma\in \sfG$ and $a\in\alg$. A natural choice is to take
$\balg=\CK(\hil)$ to be the $C^*$-algebra of compact operators on a
separable Hilbert space $\hil$, which gives a
representation $\Pim\colon \alg\to{\rm B}(\hil)$ of the algebra $\alg$ by
bounded operators ${\rm B}(\hil)$ on $\hil$ and a unitary representation
$U\colon \sfG\to{\rm B}(\hil)$ of the group $\sfG$ on~$\hil$; in this case
we call $(\Pim,U)$ a covariant representation of $(\alg,\sfG,\alpha)$ on~$\hil$.

When a group $\sfG$ acts on a space $X$, one is naturally interested
in considering the
quotient space $X/\sfG$ of $\sfG$-orbits on~$X$. When $\sfG$ acts
freely and properly on
$X$, this is described algebraically by the algebra of functions
$C_0(X/\sfG)$. More generally, the subalgebra of $\sfG$-invariant elements
$\alg^\sfG\subseteq\alg$ of a $\sfG$-algebra $\alg$ can be used to represent the quotient, even for
$\sfG$-actions with fixed points. A more general and systematic way of
dealing with the effective algebraic ``quotient'' is through the
\emph{crossed product} algebra $\alg\rtimes_\alpha \sfG$ for a dynamical
system $(\alg,\sfG,\alpha)$. For a transformation group $(X,\sfG)$,
this description is particularly powerful in the cases where the
quotient $X/\sfG$ is not a Hausdorff space, while for a free and proper
$\sfG$-action it gives an algebra with the same spectrum $X/\sfG$
as the algebra $C_0(X/\sfG)=C_0(X)^\sfG$ of
$\sfG$-invariant functions on $X$.

In order to define the crossed product
algebra, we first define
\[
\|f\|_{\rm univ} := \sup_{({\mit\Pi},U)} \big\|({\mit\Pi}\rtimes_\alpha U)(f)\big\|
\]
for compactly supported functions $f\in C_{\rm c}(\sfG,\alg) $, where the supremum is taken
over (possibly degenerate) covariant representations
$({\mit\Pi},U)$ of $(\alg, \sfG, \alpha)$ with
\[
(\Pim\rtimes_\alpha U)(f) := \int_\sfG
\Pim\big(f(\gamma)\big) U_\gamma\, \dd\mu_\sfG(\gamma) ,
\]
and $\mu_\sfG$ denotes the left invariant Haar measure on $\sfG$. This defines a norm, called the universal norm, on the
space $C_{\rm c}(\sfG,\alg)$.
Then the crossed product algebra $\alg\rtimes_\alpha \sfG$ is the
completion (in the universal norm) of the algebra $C_{\rm c}(\sfG,\alg)$
equipped with the convolution product
\begin{gather}\label{eq:convproduct}
(f\star f')(\gamma) := \int_\sfG\, f(\gamma') \alpha_{\gamma'}\big(f'\big(\gamma'^{-1} \gamma\big)\big)\, \dd \mu_\sfG(\gamma')
,
\end{gather}
for all $f,f'\colon \sfG\to\alg$. In general this is a noncommutative
multiplication, even for commutative dynamical systems.
Since $\alg$ is a $C^*$-algebra, there is a $*$-structure on the
convolution algebra defined by
\[
f^\dag(\gamma):= \Delta_\sfG(\gamma)^{-1} \alpha_\gamma\big(f\big(\gamma^{-1}\big)^*\big) ,
\]
where $\Delta_\sfG\colon \sfG\to \real^+$ is the modular function of
$\sfG$ defined through
\[
\Delta_\sfG(\gamma')
\int_\sfG\,f(\gamma)\,\dd\mu_\sfG(\gamma)=\int_\sfG f(\gamma \gamma')\,\dd\mu_\sfG(\gamma)
\]
for $f\in C_{\rm c}(\sfG,\alg)$ and $\gamma'\in \sfG$. By the uniqueness of the left
invariant Haar measure $\mu_\sfG$ up to a~positive
constant, $\Delta_\sfG(\gamma')$ is independent of~$f$
and $\Delta_\sfG$ is easily proven to be a continuous group homomorphism from~$\sfG$ to the
multiplicative group~$\real^+$; it is trivial for abelian groups and for compact
groups.

When $\alg=C_{\rm c}(X)$
is the algebra of a commutative dynamical system, the convolution
algebra $C_{\rm c}(\sfG\times X)$ consists of functions $f\colon\sfG\times X\to\complex$ and the
convolution product reads as
\[
(f\star f')(\gamma,x) = \int_\sfG f(\gamma',x)
f'\big(\gamma'^{-1} \gamma,\gamma'^{-1}\cdot x\big)\, \dd \mu_\sfG(\gamma') ,
\]
while the $*$-algebra structure is given by
\[
f^\dag(\gamma,x) = \Delta_\sfG(\gamma)^{-1} \overline{f\big(\gamma^{-1},\gamma^{-1}\cdot x\big)} .
\]

\setcounter{Theorem}{0}\setcounter{equation}{1}
% Original source lines 812--1129
\begin{Example}[noncommutative two-tori]\label{ex:NCtori}
The noncommutative torus is a fundamental example
of a noncommutative space in both physics and mathematics. Its
original incarnation~\cite{Rieffel1981} is a nice example of a crossed
product construction, which will play a fundamental role later on in this
paper. Let $\big(C(\torus),\Z,\tau^\theta\big)$ be the commutative $C^*$-dynamical system
where $\tau^\theta$ is induced through pullback by rotations of the circle $\torus$ through
a fixed angle $\theta\in\real/\Z$:
\[
\tau^\theta_n(\mathtt{f})(z) = \mathtt{f}\big(\e^{2\pi\ii n \theta} z\big) ,
\]
for $n\in\Z$, $\mathtt{f}\in C(\torus)$ and $z\in\torus$. The resulting crossed
product
\[
\A_\theta:=C(\torus)\rtimes_{\tau^\theta}\Z
\]
is a called a \emph{rotation
algebra}, and for irrational values of $\theta$ it can be identified as a noncommutative two-torus
$\torus_\theta^2$ in the following way.

By definition, the algebra
$\A_\theta$ is the universal norm completion of the convolution
algebra $C_{\rm c}(\Z\times\torus)$, whose elements $f=\{f_n\}_{n\in\Z}$ can be regarded as
sequences (with only finitely many nonvanishing
terms) of functions $f_n\colon \torus\to\complex$. The convolution product is
given by
\[
(f\star_\theta g)_n(z) := \sum_{n'\in\Z} f_{n'}(z) g_{n-n'}\big(\e^{2\pi\ii n' \theta} z\big) ,
\]
and the $*$-algebra structure is
\[
f^\dag_n(z):=\overline{f_{-n}\big(\e^{2\pi\ii n \theta} z\big)} .
\]
Via the Fourier transformation
\[
f(z,w) := \sum_{n\in\Z} f_{n}(z) w^n
\]
for $w\in\torus$, we may regard the convolution algebra $C_{\rm
 c}(\Z\times\torus)$ as a subspace of the space of functions $C\big(\torus^2\big)$ equipped
with the star-product
\begin{gather}\label{eq:starprodNCtorus}
(f\star_\theta g)(z,w) = \sum_{n\in\Z} (f\star_\theta g)_n(z) w^n .
\end{gather}
After a further Fourier transformation
\[
f_n(z) = \sum_{m\in\Z} f_{m,n} z^m
\]
and some simple redefinitions of the Fourier series involved, the
star-product \eqref{eq:starprodNCtorus} may be written in the form
\[
(f\star_\theta g)(z,w) = \sum_{(m,n)\in\Z^2}
\bigg(\sum_{(m',n')\in\Z^2} f_{m',n'} g_{m-m',n-n'}
\e^{2\pi\ii (m-m') n' \theta}\bigg) z^m w^n .
\]
This recovers the usual commutative pointwise multiplication of functions in
$C\big(\torus^2\big)$ for $\theta=0$. For $\theta\neq0$ it realizes the irrational rotation algebra $\A_\theta$ as a
deformation of the algebra of functions $C\big(\torus^2\big)$ on a
two-dimensional torus $\torus^2$; it is equivalent to the usual strict
deformation quantization of $\torus^2$ whose star-product is a twisted
convolution product on $C\big(\torus^2\big)$.

In the language of covariant representations of the dynamical system
$(C(\torus),\Z,\tau^\theta)$, the crossed product $\A_\theta$ is the
universal $C^*$-algebra generated by two unitaries $U$ and $V$
satisfying the relation~\cite[Proposition~2.56]{Williams2007}
\begin{gather}\label{eq:C*UV}
U V=\e^{-2\pi\ii \theta} V U .
\end{gather}
A concrete representation
of $\A_\theta$ on the Hilbert space $\hil = {\rm L}^2(\torus)$ is
given by defining
\[
U(\mathtt{f})(z)=z \mathtt{f}(z) \qquad \mbox{and} \qquad V(\mathtt{f})(z) =
\mathtt{f}\big(\e^{2\pi\ii\theta} z\big) .
\]
\end{Example}

\begin{Example}[{{\bf Noncommutative $\mbf{d}$-tori}}]
\label{ex:NCtorid}
The natural higher-dimensional generalization of
Example~\ref{ex:NCtori} involves a skew-symmetric real $d{\times}d$
matrix $\Theta=(\theta_{ij})$, see~\cite{Rieffel1990}. Then the noncommutative $d$-torus
$\A_\Theta=\torus_\Theta^d$ is the universal $C^*$-algebra generated
by $d$ unitaries $U_1,\dots,U_d$ satisfying the relations
\[
U_i U_j = \e^{-2\pi\ii\theta_{ij}} U_j U_i
\]
for $i,j=1,\dots,d$. By~\cite[Lemma~1.5]{Phillips2006}, every
noncommutative torus $\torus_\Theta^d$ can be obtained as an iterated
crossed product by $\Z$ in the following way. Let
$\Theta_{|d-1}=(\theta_{ij})_{1\leq i,j\leq d-1}$, and let
$U_1,\dots,U_{d-1}$ be the standard generators of
$\A_{\Theta_{|d-1}}=\torus^{d-1}_{\Theta_{|d-1}}$. Define a
group homomorphism
\smash{$\tau^{\vec\theta}\colon \Z\to\Aut(\A_{\Theta_{|d-1}})$} by
\[
\tau_n^{\vec\theta}(U_i) = \e^{2\pi\ii n \theta_{id}} U_i
 ,
\]
for $n\in\Z$, where $\vec\theta:=(\theta_{id})\in\R^{d-1}$.
Then there is an isomorphism of $C^*$-algebras
\begin{gather}\label{eq:AThetacrossed}
\A_\Theta\simeq\A_{\Theta_{|d-1}}
\rtimes_{\tau^{\vec\theta}} \Z .
\end{gather}

In the particular case where $\Theta_{|d-1}=\mbf 0_{d-1}$, we denote the
corresponding noncommutative $d$-torus by
$\A_{\vec\theta}=\torus_{\vec\theta}^d\,$, and~\eqref{eq:AThetacrossed}
shows that it can be obtained by a crossed product of the commutative
algebra of functions on a $d{-}1$-torus by an action of~$\Z$:
\[
\A_{\vec\theta} \simeq C\big(\torus^{d-1}\big)\rtimes_{\tau^{\vec\theta}}\Z .
\]
\end{Example}

\subsection{Semi-direct products and group algebras}\label{sec:semidirect}

Most of our considerations later on will focus on spaces that can be
obtained from semi-direct products of groups. We will now explain the
relation between crossed products and semi-direct products which will
be useful for these examples.

There are two ways to think about the semi-direct product construction:
\begin{enumerate}\itemsep=0pt
\item[(1)] Let $\sfG$ be a
group with two subgroups $\sfN $ and $\sfH$ such that $\sfN $ is normal. If $\sfN \cap
\sfH=\{e\}\subset \sfG$ and every
element of $\sfG$ can be written as a product of an element of
$\sfN $ with an element of $\sfH$, then we say that $\sfG$ is a semi-direct
product of its subgroups $\sfN $ and $\sfH$ and we write $\sfG=\sfN \,\sfH$.

\item[(2)]
Let $\sfN $ and $\sfH$ be two groups together with a left action
$\varphi\colon \sfH\to\Aut(\sfN )$ of $\sfH$ on $\sfN $ by automorphisms, which we denote by
$\varphi_h(n)={}^hn$ for $h\in \sfH$ and $n\in \sfN $; in particular
${}^h(n n')={}^hn\, {}^hn'$. We write ${}^\sfH \sfN $ to
indicate that $\sfH $ acts on $\sfN $ from the left. The semi-direct product
of $\sfN $ and $\sfH $ is the group $\sfN \rtimes_\varphi \sfH $ defined to be the set
$\sfN \times \sfH $ with the product
\[
(n,h)\,(n',h') = \big(n\,{}^hn',h h'\big) .
\]
The inverse is then $(n,h)^{-1}=\big({}^{h^{-1\!}}{n^{-1}},h^{-1}\big)$.
\end{enumerate}
These two definitions are equivalent: Given subgroups $\sfN ,\sfH \subset \sfG$ as
in point (1), it follows that $\sfG\simeq \sfN \rtimes_{\rm Ad} \sfH $
where ${\rm Ad}$ is the adjoint action: ${\rm Ad}_h(n)=h n h^{-1}$. On the other
hand, every element of the group $\sfG=\sfN \rtimes_\varphi \sfH $ defined in~(2) can be written as $(n,h)=(n,e_\sfH) (e_\sfN,h)$ and the subgroups $\sfN \times\{e_\sfH\}$
and $\{e_\sfN\} \times \sfH $ intersect only in the identity of $\sfG$.
If the action of $\sfH $ on $\sfN $ is trivial,
i.e., $\varphi_h={\rm id}_\sfN $ for all $h\in \sfH $, then the semi-direct
product reduces to the direct product $\sfN \rtimes_\varphi
\sfH =\sfN \times \sfH $.

Later on we will need to consider the interplay between semi-direct
products and quotient groups, which is provided by the simple
\begin{Lemma}
\label{lem:semidirectquotient}
Let $\sfG=\sfN\rtimes_\varphi\sfH$ be a semi-direct
product, and let $\V\subset \sfN$ be a subgroup which is normal in
$\sfG$. Then the quotient group $\sfG/\V$ is the
semi-direct product $(\sfN/\V)\rtimes_{\varphi^\V}\sfH$, where
${\varphi^\V}$ is the action $\varphi$ of $\sfH$ induced on the
quotient group $\sfN/\V$.
\end{Lemma}

If $\sfN $ is a group, we write $C^*(\sfN )$ for the corresponding group
$C^*$-algebra, i.e., for the crossed product $\complex\rtimes \sfN $. If
$\sfN $ is finite, then $C^*(\sfN )=\complex[\sfN ]$ is the linear space
freely generated over $\complex$ by the group elements, made into an algebra by
linearly extending the product from $\sfN $ to $\complex[\sfN ]$; equivalently
it is the algebra of continuous functions
on $\sfN $ with the convolution product.
 Given a
left $\sfH $-action $\varphi\colon \sfH \to\Aut(\sfN )$, there is an induced action
$\varphi^*\colon \sfH \to \Aut(C^*(\sfN ))$ via pullback. For $\sfH $ and $\sfN $ finite the
vector spaces $ C^*(\sfN )\times \sfH $ and $C^*(\sfN \times \sfH )$ are canonically
isomorphic, and it is straightforward
to show that the corresponding crossed product and semi-direct product are
related by
\begin{Proposition}\label{cpsp}
If $\sfN$ and $\sfH$ are finite groups, then
\[
C^*(\sfN )\rtimes_{\varphi^*}\sfH \simeq C^*(\sfN \rtimes_\varphi \sfH ) .
\]
\end{Proposition}
\begin{proof} Note that $C^*(\sfN )\rtimes_{\varphi^*}\sfH =\complex[\sfN ]\rtimes_{\varphi^*}\sfH $ is the
 vector space of functions $f\colon \sfH \to \complex[\sfN ]$ with convolution product
\[
\big(f\star_{\complex[\sfN ]\rtimes_{\varphi^*}\sfH } f'\big)(h)=\sum_{h'\in
 \sfH }\,f(h')\star_{\complex[\sfN ]}\varphi^*_{h'}\big(f'\big(h'^{-1} h\big)\big) ,
\]
and using the convolution product in $\complex[\sfN]$ this can be
written as
\[
\big(f\star_{\complex[\sfN ]\rtimes_{\varphi^*}\sfH }
f'\big)(n,h)=\sum_{h'\in \sfH} \sum_{n'\in
 \sfN } f(n',h') f'\big({}^{h'^{-1\!}}\big(n'^{-1} n\big),h'^{-1} h\big) ,
\]
which is easily seen
 to coincide with the convolution product in $ C^*(\sfN \rtimes_\varphi \sfH )=\complex[\sfN \rtimes_\varphi \sfH ]$.
\end{proof}

A more general result holds if $\sfN $ and $\sfH $ are locally compact groups with $\varphi\colon \sfH \to \Aut (\sfN )$
a~continuous action of $\sfH $ on $\sfN $ via automorphisms
(i.e., $(h,n)\mapsto \varphi_h(n)$ is a continuous map from $\sfH \times \sfN $ to $\sfN $).
In this case the semi-direct product $\sfN \rtimes_\varphi \sfH $ is a locally compact
group (in the product topology on $\sfN \times \sfH $) with $\sfN $ a closed normal subgroup and $\sfH $ a closed
subgroup (see~\cite[Proposition~3.11]{Williams2007} with
$A=\complex$). The analogue of item~(1) above also holds in the
context of locally compact groups if $\sfG$ is $\sigma$-compact, and
$\sfN$ and $\sfH$ are closed subgroups of~$\sfG$.

The action $\beta$
defining the $C^*$-dynamical system $(C^*(\sfN ),\sfH ,\beta)$ and hence the crossed
product $C^*(\sfN )\rtimes_\beta \sfH $ is the composition
of the pullback $\varphi^*$ of the action $\varphi\colon \sfH \to \Aut(\sfN )$ with
the action $\sigma_\sfH\colon \sfH \to \real^+\subset \Aut(C^*(\sfN ))$ that enters the
definition of the Haar measure on $\sfN \rtimes_\varphi \sfH $ in terms of the Haar
measures on $\sfN $ and $\sfH $: If $\mu_\sfN $ is a (left invariant)
Haar measure on $\sfN $, then the integral
$I_h(F)=\int_\sfN F\big({}^hn\big)\, \dd\mu_\sfN (n)$ for $F\in C^*(\sfN )$
is left invariant, i.e.,
$I_h(\lambda_{n'}F)= I_h(F)$ where $(\lambda_{n'}F)(n):=F\big(n'^{-1} n\big)$
for all $n'\in \sfN $
(use invariance of the Haar measure under $n\to {}^{h^{-1\!}}n'$).
Uniqueness of the Haar measure up to a positive constant then
implies there exists a function $\sigma_\sfH\colon \sfH\to\real^+$ such that
\begin{gather}\label{sigmaI}
\sigma_\sfH(h) \int_\sfN F\big({}^hn\big) \, \dd\mu_\sfN (n)=\int_\sfN F(n)
\, \dd\mu_\sfN (n) .
\end{gather}
It is straighforward to see that
$\sigma_\sfH$ is a group homomorphism and that it is
continuous~\cite[Section~2]{Williams2007}. The Haar measure
$\mu_{\sfN\rtimes_\varphi\sfH}$ on
$\sfN \rtimes_\varphi \sfH $ is then given by
\[
\int_{\sfN\rtimes_\varphi\sfH} f(n,h) \,
\dd\mu_{\sfN\rtimes_\varphi\sfH}(n,h) := \int_\sfH \int_\sfN\,
f(n,h) \sigma_\sfH(h)^{-1} \, \dd\mu_\sfN (n) \, \dd\mu_\sfH (h) .
\]
This is trivially invariant under the left $\sfN $-action, and
it is also invariant under the left $\sfH $-action $(n,h)\mapsto
(1,h') (n,h)=\big({}^{h'\!}n, h' h\big)$, using \eqref{sigmaI} with $F\big({}^{h'\!}n\big) :=f\big({}^{h'\!}n, h' h\big)$
and recalling that~$h$ is fixed in~\eqref{sigmaI}.

\begin{Theorem}\label{cpsp2}
Let $\sfN $ and $\sfH $ be locally compact groups
 and $\varphi \colon \sfH \to \Aut(\sfN )$ a continuous action of $\sfH $ on $\sfN $. Define
$\beta\colon \sfH \to \Aut(C^*(\sfN ))$ by $(\beta_{h'}\ell)
(n)=\sigma_\sfH(h')^{-1} \ell\big({}^{h'^{-1\!}}n\big)$ for all $h'\in \sfH $ and
$\ell\in C_{\rm c}(\sfN )$. Then
\[
C^*(\sfN )\rtimes_{\beta}\sfH \simeq C^*(\sfN \rtimes_\varphi
\sfH ) .
\]
\end{Theorem}

For a full
proof of Theorem~\ref{cpsp2} that takes into account the topological
and $C^*$-algebraic aspects, see~\cite[Proposition~3.11]{Williams2007}. Here
we shall just show that under the canonical injection $C_{\rm c}(\sfN \rtimes_\varphi \sfH )\hookrightarrow
C_{\rm c}(\sfN )\rtimes_\beta \sfH $, given by $f(n,h)\mapsto f(h)$ where
$f(h)(n)=f(n,h)$, the convolution product in $C_{\rm c}(\sfN \rtimes_\varphi
\sfH )$ is mapped to the convolution product in $C_{\rm c}(\sfN )\rtimes_\beta \sfH $.
Let $f,f'\in C_{\rm c}(\sfN \rtimes_\varphi \sfH )$, then
\begin{gather}\label{eq:imagefstarf'}
\big(f\star_{C_{\rm c}(\sfN \rtimes_\varphi
 \sfH )}f'\big)(n,h)=\int_\sfH \int_\sfN f(n',h')
f'\big((n',h')^{-1} (n,h)\big) \sigma_\sfH(h')^{-1} \, \dd\mu_\sfN (n') \, \dd\mu_\sfH (h').
\end{gather}
On the other hand, for the images of $f$, $f'$ in $C_{\rm c}(\sfN
)\rtimes_\beta \sfH $ we have
\[
\big(f\star_{C_{\rm c}(\sfN )\rtimes_\beta
 \sfH }f'\big)(h)=\int_\sfH f(h')\star_{C^*(\sfN
 )}\beta_{h'}\big(f'\big(h'^{-1} h\big)\big) \, \dd\mu_\sfH (h') .
\]
Using the convolution product in $C^*(\sfN)$ this can be written as
\begin{align*}
\big(f\star_{C_{\rm c}(\sfN )\rtimes_\beta
 \sfH }f'\big)(h)(n)=\int_\sfH \int_\sfN
 f(h')(n') \beta_{h'}\big(f'\big(h'^{-1} h\big)\big)\big(n'^{-1} n\big)
\, \dd\mu_\sfN (n') \, \dd\mu_\sfH (h') ,
\end{align*}
which from the definition of the action $\beta$ is easily seen to
equal the image $(f\star_{C_{\rm c}(\sfN \rtimes_\varphi \sfH
 )}f')(h)(n)$ in $C_{\rm c}(\sfN )\rtimes_\beta
 \sfH $ of the product $(f\star_{C_{\rm c}(\sfN \rtimes_\varphi
 \sfH )}f')(n,h)$ in $C_{\rm c}(\sfN \rtimes_\varphi
\sfH )$ from~\eqref{eq:imagefstarf'}.

More generally we have~\cite[Proposition~3.11]{Williams2007}
\begin{Theorem}\label{cpsp3}
Let $(\alg,\sfN \rtimes_\varphi \sfH ,\alpha)$ be a $C^*$-dynamical
system for the semi-direct product group $\sfN \rtimes_\varphi \sfH $. Then $(\alg\rtimes_{\alpha|_\sfN }\sfN ,\sfH ,\beta)$ is a $C^*$-dynamical
system, where
\[
\beta\colon \ \sfH \longrightarrow \Aut(\alg\rtimes_{\alpha|_\sfN }\sfN ) ,
\qquad h\longmapsto \beta_h
\]
is defined by
$\big(\beta_h(f)\big)(n)=\sigma_\sfH(h)^{-1} \alpha_h\big(f\big(^{h^{-1\!}}n\big)\big)$
for all $f\in C_{\rm c}(\sfN ,\alg)\subset \alg\rtimes_{\alpha|_\sfN }\sfN $, with $\sigma_\sfH\colon
\sfH \to \real^+$ defined by~\eqref{sigmaI} and $^{h^{-1\!}}n=\varphi_{h^{-1}}(n)$.
Moreover, the canonical injection $C_{\rm c}(\sfN \rtimes_\varphi \sfH ,\alg)\hookrightarrow
C_{\rm c}\big(\sfH ,C_{\rm c}(\sfN ,\alg)\big)$
extends to a~$C^*$-algebra isomorphism
\[
\alg\rtimes_{\alpha}\big(\sfN \rtimes_\varphi
\sfH \big) \simeq \big(\alg\rtimes_{\alpha|_\sfN } \sfN
\big)\rtimes_\beta \sfH .
\]
\end{Theorem}

Theorem \ref{cpsp2} is then recovered by setting $\alg=\complex$.

\setcounter{Theorem}{9}\setcounter{subsection}{3}\setcounter{equation}{8}
% Original source lines 1273--1550
\subsection{Morita equivalence and Green's theorem}\label{sec:Morita}

Crossed products of algebras provide a host of examples of dualities
which come in the form of various levels of strong and weak
equivalences of algebras, see, e.g.,~\cite{Brodzki2006}. The most
primitive form of such dualities is provided by (strong) Morita
equivalence~\cite{Rieffel1974}. A bimodule for a pair of algebras $\alg$ and $\balg$ is a
vector space $\Mcal$ which is simultaneously
 a left $\alg$-module and a right $\balg$-module, where the left action
 of $\alg$ commutes with the right action of $\balg$: $(a\cdot\xi)\cdot
 b=a\cdot (\xi\cdot b)$ for all $a\in \alg$, $b\in
 \balg$ and $\xi\in \Mcal$. If $\alg$ and $\balg$ are $C^*$-algebras,
 we say that a bimodule $\Mcal$ is an $\alg$--$\balg$ \emph{Morita
 equivalence bimodule} (or \emph{imprimitivity bimodule}) if it
is equipped with an $\alg$-valued inner product
${}_\alg\langle\,\cdot\,,\,\cdot\,\rangle$ and a $\balg$-valued inner
product $\langle\,\cdot\,,\,\cdot\,\rangle_\balg$ satisfying the
associativity condition
\[
{}_\alg\langle \psi,\phi\rangle \cdot \xi =
\psi\cdot\langle\phi,\xi\rangle_\balg ,
\]
for all $\psi,\phi,\xi\in\Mcal$, under which $\Mcal$ is complete in
the norm closures, and such that the ideal ${}_\alg\langle
\Mcal,\Mcal\rangle$ is dense in $\alg$ and $\langle
\Mcal,\Mcal\rangle_\balg$ is dense in $\balg$. The bimodule $\Mcal$ establishes a
\emph{Morita equivalence} between the algebras $\alg$ and $\balg$, and
in this case we write $\alg\ \mbf\sim_{\text{\tiny M}}\ \balg$.

Morita equivalent $C^*$-algebras have equivalent categories of
nondegenerate $*$-representations: If $\Pim_\balg\colon \balg\to{\rm B}(\hil_\balg)$ is a
representation of $\balg$ on a Hilbert space $\hil_\balg$, then we can
construct another Hilbert space
\[
\hil_\alg:=\Mcal\otimes_\balg\hil_\balg
\]
which is the quotient of the tensor product $\Mcal\otimes\hil_\balg$ by the relation $(\xi\cdot
b)\otimes\psi-\xi\otimes\Pim_\balg(b)\psi=0$ identifying the
$\balg$-actions for
$\xi\in\Mcal$, $b\in\balg$ and $\psi\in\hil_\balg$. The inner product
on $\hil_\alg$ is given by
\[
\big\langle\xi\otimes_\balg\psi\big|
\xi'\otimes_\balg\psi'\big\rangle_{\hil_\alg} :=
\big\langle\psi\big|\Pim_\balg\big(\langle\xi,\xi'\rangle_\balg
\big)\psi'\big\rangle_{\hil_\balg} ,
\]
and a representation $\Pim_\alg\colon \alg\to\rmB(\hil_\alg)$ of the algebra
$\alg$ is defined by
\[
\Pim_\alg(a)(\xi\otimes_\balg\psi) = (a\cdot\xi) \otimes_\balg\psi
\]
for $a\in\alg$ and $\xi\otimes_\balg\psi\in\hil_\alg$; this
representation is unitary equivalent to the representation
$\Pim_\balg$. Conversely, starting with a representation of $\alg$, we
can use a conjugate $\balg$--$\alg$ equivalence bimodule
$\overline{\Mcal}$ to construct a unitary equivalent representation of
$\balg$; then there are surjective bimodule homomorphisms
$\Mcal\otimes_\balg\overline{\Mcal}\to\alg$ and
$\overline{\Mcal}\otimes_\alg\Mcal\to\balg$ which satisfy a certain
transitivity law. As a particular consequence of this equivalence, Morita
equivalent algebras have homeomorphic spectra and isomorphic K-theory groups.

\begin{Example}[noncommutative two-tori]\label{ex:NCT2Morita}
A famous example of Morita equivalence in both mathematics and string
theory is provided by the
noncommutative tori ${\sf A}_\theta=\torus_\theta^2$ from
Example~\ref{ex:NCtori}. Firstly, notice from \eqref{eq:C*UV} that
changing the coset representative $\theta\in\real/\Z$ yields an identical
algebra: $\A_{\theta+m}=\A_{\theta}$ for all $m\in\Z$. Secondly, there
is an obvious
$C^*$-algebra isomorphism $\A_{-\theta}\simeq\A_\theta$ obtained by
interchanging the two generators $U$ and $V$. The converse
is also true~\cite{Rieffel1981,Rieffel1982}:
$\A_{\theta'}\simeq\A_{\theta}$ if and only if $\theta'=\theta \ {\rm
 mod}~1$. More generally, two irrational rotation $C^*$-algebras
$\A_\theta$ and $\A_{\theta'}$ are
Morita equivalent if and only if $\theta$ and $\theta'$ lie in the
same orbit under the action of $\sfGL(2,\Z)$ by fractional linear
transformations
\[
\theta'={\tt M}[\theta]:=\frac{a \theta+b}{c \theta+d} \qquad
\mbox{for} \quad {\tt M}=\bigg(\begin{matrix} a & b \\ c &
 d \end{matrix}\bigg) \in \sfGL(2,\Z) .
\]
The explicit Morita equivalence bimodules can be found
in~\cite{Rieffel1981}. On the other hand, the rational rotation
algebras $\A_\theta$ are all Morita equivalent to the commutative
algebra $C\big(\torus^2\big)$ of functions on the
two-torus~\cite{Rieffel1982}.
\end{Example}

\begin{Example}[noncommutative $d$-tori]
\label{ex:NCTdMorita}
The Morita equivalences of Example~\ref{ex:NCT2Morita} generalize to
the higher-dimensional noncommutative tori $\A_\Theta=\torus^d_\Theta$
from Example~\ref{ex:NCtorid} in the following
way~\cite{Rieffel1998}. First of all, the algebra $\A_\Theta$ is
unchanged if the matrix $\Theta$ is written in another basis of~$\Z^d$: if $B\in\sfGL(d,\Z)$ with transpose $B^{\rm t}$, then there is
a $C^*$-algebra isomorphism
$\A_{B^{\rm t}\,\Theta\,B}\simeq\A_\Theta$. More generally, consider the set of real
skew-symmetric $d{\times}d$ matrices $\Theta$ whose orbits
$\mbf M[\Theta]$ are defined for all $\mbf M\in{\sf SO}(d,d;\Z)$, where
\[
\mbf M[\Theta] = (A\,\Theta+B)\,(C\,\Theta+D)^{-1} \qquad \mbox{for} \quad
\mbf M=\left(\begin{matrix} A & B \\ C & D \end{matrix}\right) \in
{\sf SO}(d,d;\Z) ,
\]
and $A$, $B$, $C$ and $D$ are $d{\times}d$ block matrices satisfying
\[
A^{\rm t} C+C^{\rm t} A = 0 = B^{\rm t} D+D^{\rm t} B \qquad
\mbox{and} \qquad A^{\rm t} D+C^{\rm t} B=\unit_d .
\]
The set of all such matrices is dense in the space of all
skew-symmetric real $d{\times}d$ matrices, and there is a Morita
equivalence
\[
\A_{\mbf M[\Theta]} \ \mbf\sim_{\text{\tiny M}} \ \A_\Theta .
\]
However, for $d>2$ the converse is not generally true: In fact, there
are algebras $\A_\Theta$ and $\A_{\Theta'}$ that are isomorphic (and
so Morita equivalent) but for which the matrices $\Theta$ and
$\Theta'$ do not belong to the same ${\sf SO}(d,d;\Z)$
orbit~\cite{Rieffel1998}.
\end{Example}

We will also need an equivariant version of Morita equivalence in
order to show that Morita equivalent algebras induce Morita
equivalent crossed products according
to~\cite[Section~5.4]{Echterhoff2017}
\begin{Theorem}\label{equivMorita}
Let $(\alg,\sfG, \alpha)$ and $(\balg,\sfG, \beta)$ be
$C^*$-dynamical systems such that
$\alg$ and $\balg$ are Morita equivalent. Then the crossed product $C^*$-algebras
$\alg\rtimes_\alpha \sfG$ and $\balg\rtimes_\beta \sfG$ are Morita equivalent if
there exists a $\sfG$-equivariant $\alg$--$\balg$ Morita equivalence bimodule $\Mcal$,
i.e., if there is a strongly continuous action $U\colon \sfG\to \Aut(\Mcal)$ of $\sfG$ on an
$\alg$--$\balg$ Morita equivalence bimodule $\Mcal$ such that
\[
U_\gamma(a\cdot \xi)=\alpha_\gamma(a)\cdot
U_\gamma(\xi) \qquad \mbox{and} \qquad U_\gamma(\xi\cdot b)=U_\gamma(\xi)\cdot
\beta_\gamma(b) ,
\]
and
\[
{}_{\alg}\langle U_\gamma(\xi),U_\gamma(\xi')\rangle=\alpha_\gamma\big({}_\alg\langle
\xi,\xi'\rangle\big) \qquad \mbox{and} \qquad \langle
U_\gamma(\xi),U_\gamma(\xi')\rangle_\balg =\beta_\gamma\big(\langle
\xi,\xi'\rangle_\balg\big) ,
\]
for
all $\gamma\in \sfG$, $\xi,\xi'\in \Mcal$, $a\in \alg$ and $b\in \balg$.
\end{Theorem}

In this paper, our main application of Morita equivalence will involve
Green's symmetric imprimitivity theorem. Let $X$ be a locally compact space, and let
$\sfH $ and $\sfK $ be locally compact groups with commuting free and
proper actions on the right and on the left on $X$, respectively. We can lift these
actions to left actions on $C_0(X)$ by defining
$({}^h\mathtt{f})(x)=\mathtt{f}\big(h^{-1}\cdot x\big)$ and $\big({}^k\mathtt{f}\big)(x)=\mathtt{f}(x\cdot k)$
for all $\mathtt{f}\in C_0(X)$, $x\in X$, $h\in \sfH$ and $k\in \sfK $. Commutativity
of the actions of~$\sfH $ and~$\sfK $ implies that there are well-defined
induced actions of~$\sfH $ and~$\sfK $
respectively on the quotient spaces $\sfK {\setminus} X$ and $X/\sfH $, and hence
respectively on the algebras $C_0(\sfK {\setminus} X)$ and
$C_0(X/\sfH )$ which we denote~$\rt$ and~$\lt$. Green's symmetric imprimitivity theorem then reads as~\cite[Corollary~4.10]{Williams2007}
\begin{Theorem} \label{thm:Green}
There is a Morita
equivalence of $C^*$-algebras
\begin{gather}\label{eq:Greenimprim}
C_0(\sfK {\setminus} X)\rtimes_\rt \sfH \ \mbf\sim_{\text{\rm \tiny M}} \ C_0(X/\sfH )\rtimes_\lt \sfK
\end{gather}
implemented by the Morita equivalence $($or imprimitivity$)$
bimodule $\Mcal$ which is the completion of $C_{\rm c}(X)$ with the
actions
\begin{gather*}
(a\cdot\xi)(x) =\int_\sfK\,a(k,x\cdot\sfH) \xi\big(k^{-1}\cdot x\big)
\Delta_\sfK(k)^{1/2} \, \dd\mu_\sfK(k) , \\
(\xi\cdot b)(x) =\int_\sfH \xi\big(x\cdot h^{-1}\big) b\big(h,\sfK\cdot x\cdot h^{-1}\big) \Delta_\sfH(h)^{-1/2} \, \dd\mu_\sfH(h) ,
\end{gather*}
for all $x\in X$, $a\in C_{\rm c}(\sfK\times X/\sfH)$, $b\in C_{\rm
 c}(\sfH\times \sfK {\setminus} X)$ and $\xi\in C_{\rm
 c}(X)$, and the inner products
\begin{gather*}
{}_\alg\langle\xi,\xi'\rangle(k,x\cdot\sfH) = \Delta_\sfK(k)^{-1/2}
\int_\sfH \xi(x\cdot h) \overline{\xi'\big(k^{-1}\cdot x\cdot h\big)} \,
\dd\mu_\sfH(h) , \nonumber\\
\langle\xi,\xi'\rangle_\balg(h,\sfK\cdot x) =
\Delta_\sfH(h)^{-1/2} \int_\sfK \overline{\xi\big(k^{-1}\cdot x\big)}
\xi'\big(k^{-1}\cdot x\cdot h\big) \, \dd\mu_\sfK(k)  ,
\end{gather*}
for all $x\in X$, $h\in\sfH$, $k\in\sfK$ and $\xi,\xi'\in C_{\rm c}(X)$.
\end{Theorem}

Theorem~\ref{thm:Green} has several useful applications and
corollaries, see, e.g.,~\cite{Rieffel1976}. A particularly relevant
special case that we shall use below is when $\sfK $ is the trivial group,
in which case~\eqref{eq:Greenimprim} reduces to the Morita equivalence
\begin{gather*}%\label{eq:GreenKe}
C_0(X)\rtimes_\rt \sfH \ \mbf\sim_{\text{\tiny M}} \ C_0(X/\sfH ) ,
\end{gather*}
illustrating the use of crossed products in describing
quotients. This equivalence can in fact be strengthened to a stable
isomorphism~\cite{Rieffel1976}
\begin{gather}\label{eq:XsfHiso}
C_0(X)\rtimes_\rt \sfH \simeq C_0(X/\sfH ) \otimes \CK\big(\Leb^2(\sfH)\big) ,
\end{gather}
where $\CK$ denotes the algebra of
compact operators.
\begin{Example}[tori]\label{ex:Moritatori}
A particularly relevant example for us is the case $X=\real^d$ with
$\sfH=\zed^d$ acting by translations $(n,x)\mapsto x+n$ for
$n\in\zed^d$ and $x\in\real^d$, which realizes the
$d$-dimensional torus $\torus^d=\real^d/\zed^d$ as a crossed product:
\[
C\big(\torus^d\big) \ \mbf\sim_{\text{\tiny M}} \ C_0\big(\real^d\big)\rtimes_\rt \zed^d .
\]
Let us illustrate the construction explicitly. The convolution algebra
$C_{\rm c}\big(\Z^d\times\real^d\big)\!\subset\! C_0\big(\R^d\big)\rtimes_\rt\Z^d$ can be
identified with the space of sequences $f=\{f_n\}_{n\in\Z^d}$ of
functions $f_n\colon \real^d\to\complex$ with the convolution product
\[
(f\star g)_n(x) = \sum_{m\in\Z^d} f_m(x) g_{n-m}(x-m) .
\]

Consider the algebra
\[
\alg:=\big\{f\in C_0\big(\real^d,\CK\big(\ell^2\big(\zed^d\big)\big)\big) \, \big| \, f(x+m)
= U_m f(x) U_{m}^{-1}\big\} ,
\]
where $U_m$ is the unitary shift operator on $\ell^2\big(\zed^d\big)$ defined by
$(U_ma)_n=a_{n-m}$ for each $m\in\Z^d$ and $a=\{a_n\}_{n\in\Z^d}$.
Define a map ${\mit\Phi}\colon C_{\rm c}\big(\Z^d\times\real^d\big)\to
C_{\rm c}\big(\real^d,\CK\big(\ell^2\big(\Z^d\big)\big)\big)$ by
\[
\big({\mit\Phi}(f)(x)\big)_{mn} = f_{m+n}(x+n) .
\]
It is easy to see that ${\mit\Phi}(f)(x+m) =
U_m {\mit\Phi}(f)(x) U_m^{-1}$ for all $f\in C_{\rm
 c}\big(\zed^d\times\real^d\big)$, and if $a=(a_{mn})\in\alg$ then defining
$f_n:=a_{0n}$ gives ${\mit\Phi}(f)=a$. It is also easy to check that
${\mit\Phi}(f\star g) = {\mit\Phi}(f) {\mit\Phi}(g)$, and consequently
$\mit\Phi$ gives an algebra isomorphism
\[
{\mit\Phi}\colon \ C_0\big(\real^d\big)\rtimes_\rt\zed^d \xrightarrow{ \ \simeq \ } \alg .
\]

The explicit Morita equivalence bimodule is now obtained from the
completion of
\[
\Mcal = \big\{\xi\in C_{\rm c}\big(\real^d,\ell^2\big(\Z^d\big)\big) \, \big| \, \xi(x+m) =
U_m \xi(x) \big\} .
\]
The left action of the algebra $\alg={\mit\Phi}\big(C_{\rm c}\big(\Z^d\times\real^d\big)\big)$ is by left
matrix multiplication on $\Mcal$:
\begin{align*}
\alg\times\Mcal &\longrightarrow\Mcal , \\
(a,\xi) & \longmapsto a\cdot\xi , \qquad (a\cdot\xi)_n=\sum_{m\in\Z^d}
a_{nm} \xi_m ,
\end{align*}
while the right action of the algebra $C\big(\torus^d\big)$ is by right pointwise
multiplication on~$\Mcal$:
\begin{align*}
\Mcal\times C\big(\torus^d\big) & \longrightarrow\Mcal , \nonumber \\
(\xi,b) & \longmapsto \xi\cdot b , \qquad (\xi\cdot b)_n=\xi_n b .
\end{align*}
The left and right inner
products are respectively given by
\begin{gather*}
{}_\alg\langle\xi,\eta\rangle = \xi\otimes\eta^* , \\
\langle\xi,\eta\rangle_{C(\torus^d)} = \sum_{n\in\zed^d}
\overline{\xi_n} \eta_n ,
\end{gather*}
for $\xi,\eta\in\Mcal$.
Together with the isomorphism $\mit\Phi$, this establishes a
Morita equivalence between the algebras $C_0\big(\real^d\big)\rtimes_\rt\zed^d$
and $C\big(\torus^d\big)$.
\end{Example}


% Original source lines 1571--1740
\section{Topological T-duality and twisted tori}\label{sec:TopTtorus}

In this section we shall apply the results of Section~\ref{sec:NAdef}, and in particular
Green's theorem, to present a scheme
that will be employed in our study of topological T-duality. We shall
then illustrate how our scheme works to reproduce some standard
(commutative) examples of T-dual spaces.

\subsection{Twisted tori and their T-duals}\label{sec:twistedtori}

We are interested in formulating a notion of T-duality for
``torus bundles without $H$-flux'', which for our purposes can be
characterised by the following general class of spaces.

\begin{Definition}[twisted tori]
Let $\sfG$ be a locally compact group which admits a
cocompact discrete subgroup $\Lambda_\sfG$, i.e., a lattice in~$\sfG$, which we let act on $\sfG$ by left
multiplication. The quotient space
\[
\torus_{\Lambda_\sfG}:=\Lambda_\sfG\backslash \sfG
\]
is a \emph{twisted torus}.
\end{Definition}

By~\cite[Lemma~6.2]{Milnor1976}, only unimodular groups can contain lattices, i.e., groups $\sfG$ whose
modular function $\Delta_\sfG$ is identically equal to~$1$.

\begin{Example}[tori]\label{ex:tori}
Every lattice in the abelian Lie group $\sfG=\real^d$ is
isomorphic to $\Lambda_\sfG=\zed^d$, acting by translations. Then
$\torus_{\zed^d}=\zed^d\backslash\real^d=:\torus^d$ is the $d$-dimensional torus.
\end{Example}

\begin{Example}[nilmanifolds]\label{ex:nilmanifolds}
Generalizing Example~\ref{ex:tori}, let $\sfG$ be a connected and simply-connected nilpotent Lie group. Then a theorem of
Malcev~\cite[Theorem~2.12]{Raghunathan1972} establishes the existence
of a lattice $\Lambda_\sfG$ in $\sfG$ if and only if
$\sfG$ can be defined over the rationals,
i.e., there exists a basis for its Lie algebra which has rational
structure constants, and in this case $\torus_{\Lambda_\sfG}=\Lambda_\sfG\backslash \sfG$ is a nilmanifold.
\end{Example}

\begin{Example}[orbifolds]
Let $\sfG$ be a compact Lie group. Then the lattices in $\sfG$ are
precisely the finite subgroups $\Gamma$ of $\sfG$, and
$\torus_\Gamma=\Gamma\backslash\sfG$ is a smooth orbifold. For instance,
for $\sfG=\SU(2)$ the twisted tori are precisely the three-dimensional
ADE orbifolds $\torus_\Gamma=\Gamma\backslash\Sphere^3$
of the three-sphere for a finite subgroup $\Gamma\subset\SU(2)$;
for $\Gamma=\zed_n$ a cyclic subgroup of order $n\geq2$, this recovers
the familiar lens spaces $\torus_{\zed_n}=\zed_n\backslash\Sphere^3=:\mathbb{L}(n,1)$.
\end{Example}

Following~\cite{Mathai2004}, we come now to a central concept of this paper.

\begin{Definition}[topological T-duality]\label{def:Tduality}
Let $\torus_{\Lambda_\sfG}$ be a twisted torus which admits a
non-trivial right action
of the abelian Lie group $\real^n$ for some $n\geq1$. The crossed product
\[
C(\torus_{\Lambda_\sfG})\rtimes_\rt \real^n
\]
is a \emph{$C^*$-algebraic T-dual} of the twisted torus.
\end{Definition}
If the spectrum of the crossed product algebra $
C(\torus_{\Lambda_\sfG})\rtimes_\rt \real^n
$ is a Hausdorff topological space~$X$ (for instance
if it is Morita
equivalent to a commutative $C^*$-algebra $C(X)$), then we say that~$X$ is T-dual to the
twisted torus $\torus_{\Lambda_\sfG}$ and call $X$ a `classical T-dual'; otherwise we say that the T-dual of
$\torus_{\Lambda_\sfG}$ is a noncommutative space.

For Definition \ref{def:Tduality} to be a `good' notion of T-duality, we should first explain
\begin{itemize}\itemsep=0pt
\item[(a)] in
what precise sense $\torus_{\Lambda_\sfG}$ and
$C(\torus_{\Lambda_\sfG})\rtimes_\rt\real^n$ are `equivalent', and
\item[(b)] how
T-duality applied twice returns the original twisted torus.
\end{itemize}
The answers to both of these points turns out to be provided by
working in a suitable category tailored to our treatment of
topological T-duality.

\subsection[T-duality in the category $\CCK\hspace{-1mm}\CCK$]{T-duality in the category $\CCK\hspace{-1mm}\CCK$}
\label{sec:TdualityKK}

The terminology `topological T-duality' refers to a coarse equivalence at
the level of topology; for $C^*$-algebras the topology is measured by
K-theory. A more powerful refinement is provided by Kasparov's bivariant
K-theory which constructs groups ${\rm KK}(\alg,\balg)$ for any pair
of separable $C^*$-algebras $\alg$ and $\balg$; when $\alg=\complex$,
the group ${\rm KK}(\complex,\balg)\simeq {\rm
 K}(\balg)$ is the K-theory group of $\balg$. The cycles in
Kasparov's groups ${\rm KK}(\alg,\balg)$, called \emph{Kasparaov
 $\alg$--$\balg$ bimodules}, are triples $(\hil,\phi,T)$ where $\hil$
is a right Hilbert $\balg$-module, $\phi$ is a $*$-representation of $\alg$
on $\hil$, and $T\in\End_\balg(\hil)$ is a $\balg$-linear operator on
$\hil$, subject to certain compactness conditions; we do not
provide further details of the definition here and instead refer to~\cite{Brodzki2006} for a
concise review of KK-theory in the context that we shall use it
in this paper. Kasparov bimodules may be thought of as generalizations
of morphisms between $C^*$-algebras, in the sense that any algebra homomorphism $\phi\colon \alg\to\balg$
determines a class $[\phi]\in{\rm KK}(\alg,\balg)$, represented by the
$\alg$--$\balg$-bimodule $(\balg,\phi,0)$.

A key feature of Kasparov's KK-theory is the composition product
\[
\otimes_\balg\colon \ {\rm KK}(\alg,\balg)\times{\rm
 KK}(\balg,\Ccal)\longrightarrow{\rm KK}(\alg,\Ccal) ,
\]
which is bilinear and associative. This product is compatible with the
composition of morphisms $\phi\colon \alg\to\balg$ and $\psi\colon \balg\to\Ccal$
of $C^*$-algebras: $[\phi]\otimes_\balg[\psi] = [\psi\circ\phi]$. It
also makes ${\rm KK}(\alg,\alg)$ into a~ring with unit
$1_\alg=[\Id_\alg]$. We say that an element $\alpha\in{\rm
 KK}(\alg,\balg)$ is \emph{invertible} if there exists an element
$\beta\in{\rm KK}(\balg,\alg)$ such that
$\alpha\otimes_\balg\beta=1_\alg$ and
$\beta\otimes_\alg\alpha=1_\balg$.

An important special instance of Kasparov
bimodules comes from Morita equivalence: Any Morita equivalence
$\alg$--$\balg$ bimodule $\Mcal$ is also a Kasparov bimodule
$(\Mcal,\phi,0)$, with $\phi\colon \alg\to\End(\Mcal)$ the left action of
$\alg$, which defines an invertible class $[\Mcal]\in{\rm
 KK}(\alg,\balg)$ with inverse $[\,\overline{\Mcal}\,]\in{\rm
 KK}(\balg,\alg)$ given by the conjugate $\balg$--$\alg$ bimodule
$\overline{\Mcal}$. Generally, if there exists an invertible element
$\alpha\in{\rm KK}(\alg,\balg)$, then the algebras $\alg$ and $\balg$
are said to be \emph{{KK}-equivalent}, and we write $\alg\
\mbf\sim_{\text{\tiny{KK}}}\ \balg$. Thus Morita equivalence implies
KK-equivalence, but the converse is not generally true. KK-equivalent algebras have isomorphic
K-theory groups, but not necessarily homeomorphic spectra.

This refinement naturally
suggests an approach to T-duality where the category of separable
$C^*$-algebras with $*$-homomorphisms is replaced with an additive
category ${\CCK\hspace{-1mm}\CCK}$, whose objects are again separable $C^*$-algebras
but whose morphisms between any two objects $\alg$ and $\balg$ are
given by the classes in ${\rm KK}(\alg,\balg)$ (see,
e.g.,~\cite{Brodzki2007}). The composition product defines the
composition law, and isomorphic algebras
in this category are precisely the KK-equivalent
algebras; in particular, Morita equivalent algebras are isomorphic
as objects in ${\CCK\hspace{-1mm}\CCK}$. Our formulation and computations of topological
T-duality will always take place in this category, and in this
setting we can easily provide answers to points~(a) and~(b) below Definition~\ref{def:Tduality}
through
\begin{Theorem}\label{thm:equivalences}
If $\torus_{\Lambda_\sfG}$ is a twisted torus with a
non-trivial right action
of $\real^n$, then there are
isomorphisms in the category ${\CCK\hspace{-1mm}\CCK}$ given by the equivalences
\begin{itemize}\itemsep=0pt
\item[{\rm (a)}]  $C(\torus_{\Lambda_\sfG}) \
 \mbf\sim_{\text{\tiny{KK}}} \
 C(\torus_{\Lambda_\sfG})\rtimes_\rt\real^n  $ $($up to a shift of degree
 $n \ {\rm mod}~2)$, and
\item[{\rm (b)}]  $\big(C(\torus_{\Lambda_\sfG})\rtimes_\rt
 \real^n\big)\rtimes_{\widehat\rt} \real^n \ \mbf\sim_{\text{\tiny
 M}} \ C(\torus_{\Lambda_\sfG})$.
\end{itemize}
\end{Theorem}
\begin{proof}
The KK-equivalence (a) follows from the Connes--Thom isomorphism,
formulated in the language of KK-theory~\cite{Fack1981}. The Morita
equivalence~(b) follows from Takai duality
(Theorem~\ref{thm:Takaiduality}).
\end{proof}

\setcounter{equation}{2}
% Original source lines 1879--1887
More generally, suppose that the $\real^n$-action on $\torus_{\Lambda_\sfG}$ is
induced by a free and proper right action of $\real^n$ on the covering group $\sfG$
which commutes with the left action of the lattice $\Lambda_\sfG$ on~$\sfG$. We can then apply Green's theorem
(Theorem~\ref{thm:Green}) to get the Morita
equivalence
\begin{gather}\label{scheme}
C(\torus_{\Lambda_\sfG})\rtimes_\rt \real^n \ \mbf\sim_{\text{\tiny M}} \
C_0\big(\sfG/\real^n\big)\rtimes_\lt \Lambda_\sfG .
\end{gather}

\setcounter{section}{3}
% Original source lines 2091--2172
\section{Topological T-duality for almost abelian solvmanifolds}\label{sec:TdualityMostow}

A large class of twisted tori of interest as string compactifications
come in the form of fibrations over tori. These are the solvmanifolds
which are based on solvable groups $\sfG$ and generalize the
nilmanifolds discussed in Example~\ref{ex:nilmanifolds}. The fibrations
underlying these twisted tori are called {Mostow
 bundles}~\cite{Mostow1954}, and we are particularly interested in
the cases where the Mostow bundle is a torus bundle. A good source for
the material used in this section is~\cite{Bock2009} (see also~\cite{Console2016, OpreaTralle}).

\subsection{Mostow bundles}

Let $\sfG$ be
a connected and simply-connected solvable Lie group. Recall that its
nilradical $\sfN$ is the maximal connected nilpotent normal subgroup. It
has dimension $\dim\sfN\geq\frac12\dim\sfG$.

We first consider the case $\dim\sfN=\dim\sfG$. Then
$\sfN=\sfG$ and the group $\sfG$ is nilpotent. In this case, under the
conditions discussed in
Example~\ref{ex:nilmanifolds}, there exists a lattice
$\Lambda_\sfG\subset\sfG$ and the twisted torus $\torus_{\Lambda_\sfG}$ is a
nilmanifold. If $\sfG$ is abelian then $\torus_{\Lambda_\sfG}$ is a
torus. If $\sfG$ is non-abelian then there is a group extension
\[
1\longrightarrow [\sfG,\sfG]\longrightarrow\sfG\xrightarrow{ \ \pi \
}\sfG_{\rm ab}\longrightarrow 1
\]
of its commutator subgroup $[\sfG,\sfG]$,
and both $\Lambda_\sfG\cap[\sfG,\sfG]$ and $\pi(\Lambda_\sfG)$ are
lattices in the nilpotent Lie group $[\sfG,\sfG]$ and the
abelianization $\sfG_{\rm ab}:=[\sfG,\sfG]\backslash\sfG$ of $\sfG$, respectively~\cite{Corwin1990}. This exhibits
the twisted torus $\torus_{\Lambda_\sfG}=\Lambda_\sfG\backslash \sfG $ as a fibration over the torus
$\pi(\Lambda_\sfG)\backslash\sfG_{\rm ab}$ with
nilmanifold fibres,{\samepage
\[
\big(\Lambda_\sfG\cap[\sfG,\sfG]\big)\backslash[\sfG,\sfG] \longrightarrow \torus_{\Lambda_\sfG} \longrightarrow
\pi(\Lambda_\sfG)\backslash \sfG_{\rm ab} .
\]
If $[\sfG,\sfG]$ is an abelian Lie group then the twisted
torus is a torus bundle over a torus.}

Suppose now that the group $\sfG$ is not nilpotent. Then $\sfN\backslash\sfG$
is a non-trivial abelian Lie group. If $\sfG$ admits a lattice $\Lambda_\sfG$, then
$\Lambda_{\sfN}:=\Lambda_\sfG\cap\sfN$ is a lattice in $\sfN$ and
$\Lambda_\sfG \sfN=\sfN \Lambda_\sfG$ is a closed subgroup of $\sfG$, so
$\Lambda_\sfG \sfN\backslash\sfG$ is a torus. The twisted torus
$\torus_{\Lambda_\sfG}=\Lambda_\sfG\backslash\sfG$ is then a~fibration over this torus with fibre
the nilmanifold $\Lambda_{\sfN}\backslash\sfN=\Lambda_\sfG\backslash \Lambda_\sfG\,\sfN$. This
bundle is called the \emph{Mostow bundle}~\cite{Mostow1954}. We summarise these statements as
\begin{Theorem}[Mostow bundles]
\label{thm:Mostowbundle}
Let ${\Lambda_\sfG}$ be a lattice in a connected and simply-connected
solvable Lie group $\sfG$ and $\torus_{\Lambda_\sfG} =
\Lambda_\sfG\backslash\sfG$ the associated solvmanifold. Let $\sfN$ be the nilradical
 of $\sfG$.
Then $\Lambda_\sfG\,\sfN$ is a closed subgroup of $\sfG$,
$\Lambda_{\sfN}:=\Lambda_\sfG\cap\sfN$ is a lattice in
$\sfN$, and $\Lambda_\sfG \sfN\backslash\sfG$ is a torus.
It follows that the twisted torus
$\torus_{\Lambda_\sfG}$ is a fibration over this torus with
nilmanifold fibre:
\[
\Lambda_{\sfN}\backslash\sfN=\Lambda_\sfG\backslash\Lambda_\sfG \sfN
\longrightarrow
\torus_{\Lambda_\sfG}\longrightarrow \Lambda_\sfG \sfN\backslash\sfG .
\]
\end{Theorem}

\begin{Remark}\label{rem:Mostow}
The structure group of the Mostow bundle is
$\Lambda_{\sfG_0}\backslash\Lambda_\sfG \sfN$, where $\Lambda_{\sfG_0}$ is
the largest subgroup of $\Lambda_{\sfG}$ which is normal in
$\Lambda_\sfG \sfN$ (cf.~\cite{Bock2009}).
In particular, if $\Lambda_{\sfG}=\Lambda_{\sfG_0}$ then the Mostow bundle is
a principal $\Lambda_{\sfG}\backslash\Lambda_\sfG \sfN$-bundle. In
this case there is a well-defined left $\Lambda_\sfG \sfN$-action on $\torus_{\Lambda_\sfG} =
\Lambda_\sfG\backslash\sfG$ and each point has isotropy subgroup ${\Lambda_\sfG}$,
so that the induced $\Lambda_{\sfG}\backslash\Lambda_\sfG \sfN$-action
is principal.
\end{Remark}

\setcounter{Theorem}{4}\setcounter{subsection}{1}\setcounter{equation}{1}
% Original source lines 2272--2486
\subsection{Almost abelian solvmanifolds}\label{sec:abeliansolv}

In contrast to the case of nilpotent groups, there is no
simple criterion for the existence of a~lattice in a~general
connected and simply-connected solvable Lie group, as is required to define a corresponding twisted
torus. To formulate such a criterion, we specialise to \emph{almost
abelian solvable groups}: these are the solvable Lie groups $\sfG$ of
dimension $d$
whose nilradical $\sfN$ has codimension one and is abelian:
\[
\sfN \simeq \real^{d-1} .
\]
Then $\sfG$ has the
structure of a semi-direct product
\[
\sfG = \sfN\rtimes_\varphi\real
\]
for a continuous one-parameter left group action
$\varphi\colon \real\to\Aut(\sfN)$; concretely, $\varphi$ is given by the
adjoint action of the one-dimensional subgroup $\sfH=\real$ on $\sfN$
in the group $\sfG$ (cf.~Section~\ref{sec:semidirect}). This exhibits
$\sfG$ as a nontrivial group extension
\[
1\longrightarrow \sfN \longrightarrow \sfG \longrightarrow \R \longrightarrow 1 .
\]

We can regard the one-parameter group action $\varphi$ as a matrix
$\varphi_x\in\sfGL(d-1,\real)$ for each $x\in\real$ acting on the
vector space $\sfN$, which we identify with $\real^{d-1}$ via a choice
of basis $\vec e_1,\ldots,\vec e_{d-1}$. Since
$\varphi_{0}=\unit_{d-1}$ and $\varphi_x$ is always non-singular, it
follows from the continuity of $\varphi$ and the determinant that
$\det\varphi_x>0$ for all $x\in\real$. Then $\sfG$ admits a lattice
$\Lambda_\sfG$ if and
only if there exists $x_0\in\real^\times$ such that $\varphi_{x_0}$ is
conjugate to an integer matrix $\ttM\in\sfSL(d-1,\zed)$:
\begin{gather}\label{eq:monodromydef}
\Sigma^{-1} \varphi_{x_0} \Sigma=\ttM
\end{gather}
for some
$\Sigma\in\sfGL(d-1,\real)$.
In this case the twisted torus
$\torus_{\Lambda_\sfG}=\Lambda_{\sfG}\backslash \sfG$ is called an {\emph{almost abelian solvmanifold}}.
The condition~\eqref{eq:monodromydef} strongly restricts the homomorphisms
$\varphi\colon \real\to\Aut(\sfN)$; in particular, it requires that the
characteristic polynomial of $\varphi_{x_0}$ has integer
coefficients. In this case the lattice (in the standard basis $\vec
e_1,\ldots,\vec e_{d-1}$) is given by
\[%\label{standardbasis}
\Lambda_\sfG = \Sigma\cdot\zed^{d-1}\rtimes_{\varphi|_{x_0\,\zed}}
x_0 \zed ,
\]
which correspondingly sits as a nontrivial group extension
\[
1\longrightarrow \Sigma \cdot \zed^{d-1} \longrightarrow \Lambda_\sfG
\longrightarrow x_0 \zed \longrightarrow 1 .
\]
Then $\Lambda_{\sfN}\setminus\sfN\simeq\torus^{d-1}$, and the corresponding Mostow bundle realises the twisted torus
$\torus_{\Lambda_\sfG}$ as a torus bundle over a circle
$\Lambda_\sfG \sfN\backslash\sfG \simeq \torus$, whose monodromy
is specified by the matrix $\ttM$ in the mapping class group $\sfSL(d-1,\zed)$ of
orientation-preserving
automorphisms up to homotopy of the torus fibres~$\torus^{d-1}$.

For an almost abelian solvmanifold we can make the twisted torus construction more concrete
by choosing the global
coordinates $(z,x)\in\real^{d-1}\times\real$ on the group
manifold (associated with the basis $\vec e_1,\ldots,\vec e_{d-1}$).
The group multiplication of the semi-direct product
$\sfG=\real^{d-1}\rtimes_\varphi\real$ is
\begin{gather}\label{eq:grouplaw}
(z,x) (z',x') = (z + \varphi_x\cdot z',x+x') ,
\end{gather}
where we used $\varphi_x \varphi_{x'} = \varphi_{x+x'}$, and the
inverse of a group element is
\[
(z,x)^{-1} = (-\varphi_{-x}\cdot z,-x) ,
\]
where we used $\varphi_x^{-1}=\varphi_{-x}$.
The twisted torus is defined as the quotient $\torus_{\Lambda_\sfG}=\Lambda_\sfG
{\setminus} \sfG$ which is generated by the equivalence relation $(z,x)\sim
(\Sigma\cdot\gamma,x_0\,\alpha)\,(z,x)$ for all $(\gamma,\alpha)\in\zed^{d-1}\times\zed$. The global structure
of $\torus_{\Lambda_{\sfG}}$ is generated by the simultaneous local coordinate
identifications under the action of the elements
$(\Sigma\cdot\gamma,x_0 \alpha) $ of $\Lambda_\sfG$ given by
\begin{gather}
(z,x) \longmapsto (
 z+\Sigma\cdot\gamma,x) , \nonumber \\
(z,x) \longmapsto (
\Sigma {\tt M}^\alpha\,\Sigma^{-1}\cdot z,x+x_0\,\alpha) .\label{eq:Mostowtorus}
\end{gather}
These identifications explicitly exhibit the twisted torus as a torus bundle over a circle, with
local fiber coordinates $z\in \torus^{d-1}$ and base coordinate $x\in \torus$,
 whose monodromy is specified by the matrix ${\tt M}\in
 \sfSL(d-1,\zed)$, and whose periods are given respectively by
 $\Sigma\in\sfGL(d-1,\real)$ and $x_0\in\real^\times$.

\subsection{$C^*$-algebra bundles}\label{sec:C*bundles}

Mostow bundles and their $C^*$-algebraic T-duals can be grouped together
under the general heading of `$C^*$-algebra bundles', which encompasses
the notions of $C_0(X)$-algebras and $C^*$-bundles, as we now
explain; see Section~8.1 and Appendix~C of~\cite{Williams2007} for
further details.
Let~$X$ be a locally compact Hausdorff space. There are two equivalent
notions for the $C^*$-algebra analogue of a fibre bundle over $X$.

A \emph{$C_0(X)$-algebra} is a $C^*$-algebra $\alg$
equipped with a nondegenerate injection $\iota$ of $C_0(X)$
into the centre of its multiplier algebra, called the \emph{structure
 map}. For $\mathtt{f}\in C_0(X)$ and
$a\in\alg$, we abbreviate $\iota(\mathtt{f})a$ by $\mathtt{f}\cdot a$. This endows
$\alg$ with a $C_0(X)$-bimodule
structure.

Given a family $\balg=(\balg_x)_{x\in X}$ of $C^*$-algebras, a
\emph{section} of $\balg$ is
a map $s\colon X\to\balg$ such that $s(x)\in\balg_x$ for all $x\in X$; we
denote the space of sections of $\balg$ which vanish at infinity by~$\Gamma_0(\balg)$. The family $\balg$ is then called a~\emph{$C^*$-bundle} over $X$ with fibres $\balg_x$ if the following
conditions are satisfied:
\begin{itemize}\itemsep=0pt
\item
$\Gamma_0(\balg)$ is a $C^*$-algebra under pointwise operations and the
 supremum norm;
\item $\balg_x=\{s(x)\, |\, s\in\Gamma_0(\balg)\}$ for each
 $x\in X$;
\item $\Gamma_0(\balg)$ is closed under multiplication by
 $C_0(X)$; and
\item For each $s\in\Gamma_0(\balg)$, the function $x\mapsto \|s(x)\|$
is upper semi-continuous, i.e., the set
$\{x\in X \,|\, \|s(x)\|< \epsilon\}$ is open in $X$ for all $\epsilon>0$.
\end{itemize}
In this
paper we will only be concerned with $C^*$-bundles that have non-zero fibres.

If $\balg$ is a $C^*$-bundle over $X$, then its section algebra
$\Gamma_0(\balg)$ is a $C_0(X)$-algebra: its structure map $\iota$ from
$C_0(X)$ is defined by $\iota(\mathtt{f})s=\mathtt{f} s$. Conversely, if $\alg$ is a
$C_0(X)$-algebra, then the fibre~$\alg_x$ of $\alg$ over $x\in X$ is $\alg_x:=\alg/\CI_x$, where
$\CI_x=\{\mathtt{f}\cdot a\, |\, \mathtt{f}\in C_0(X) , \,
\mathtt{f}(x)=0 , \, a\in\alg\}$ is identified as the ideal in
$\alg$ of sections vanishing at $x$. If
$a\in\alg$, we write $a(x)=a+\CI_x$ for its image in~$\alg_x$. The
function $x\mapsto\|a(x)\|$ is upper semi-continuous
and vanishes at infinity with
\[
\|a\| = \sup_{x\in X} \big\|a(x)\big\|
\]
for all $a\in\alg$. The
elements $a\in \alg$ can in this way be viewed as sections of a
$C^*$-bundle $(\alg_x)_{x\in X}$. We will sometimes use the
notation $\coprod_{x\in X} \alg_x$ for $\alg$ when we wish to emphasise
its structure as a $C^*$-bundle over $X$ with fibre $C^*$-algebras
$\alg_x$. These definitions do not require
local triviality of the bundle nor the fibres of the bundle to be
isomorphic to one another. $C^*$-algebra bundles over $X$ are objects of a
category whose morphisms are \emph{fibrewise} $*$-homomorphisms,
i.e., $C_0(X)$-linear morphisms $\psi\colon \alg\to\balg$: $\psi(\mathtt{f}\cdot
a)=\mathtt{f}\cdot\psi(a)$ for all $\mathtt{f}\in C_0(X)$ and $a\in\alg$; then $\psi$
induces $*$-homomorphisms $\psi_x\colon \alg_x\to\balg_x$ such that
$\psi_x\big(a(x)\big)=\psi(a)(x)$ for all $a\in\alg$.

\begin{Example}[trivial $C^*$-algebra bundles]
If $\Dcal$ is any $C^*$-algebra, then $\alg=C_0(X,\Dcal)\simeq
C_0(X)\otimes\Dcal$ is naturally a $C_0(X)$-algebra with structure map
\[
\big(\iota(\mathtt{f})a\big)(x) := \mathtt{f}(x)\,a(x)
\]
for $\mathtt{f}\in C_0(X)$, $a\in\alg$ and $x\in X$. In this case each fibre
$\alg_x$ is canonically identified with~$\Dcal$ and elements of~$\alg$
are obviously identified with sections.
\end{Example}

\begin{Example}[continuous maps]\label{ex:ctsmaps}
Let~$X$ and~$Y$ be locally compact spaces and $\sigma\colon Y\to X$ a~continuous surjective map. Then $C_0(Y)$ is a $C_0(X)$-algebra with structure map
$\iota(\mathtt{f})\mathtt{g}:=(\mathtt{f}\circ\sigma)\,\mathtt{g}$, for $\mathtt{f}\in C_0(X)$ and $\mathtt{g}\in C_0(Y)$,
and fibers $C_0(Y)_x\simeq C_0\big(\sigma^{-1}(x)\big)$.
\end{Example}

In this paper we are particularly interested in crossed
products of $C^*$-algebra bundles. Let~$\alg$ be a $C_0(X)$-algebra,
and denote by $\Aut_X(\alg)$ the group of {fibrewise
automorphisms} of~$\alg$. A~\emph{fibrewise action} of a locally compact group $\sfG$ on
$\alg$ is then a group homomorphism $\alpha\colon \sfG\to\Aut_X(\alg)$. This
implies that $\alpha$ induces an action~$\alpha^x$ on each fiber
$\alg_x$ for $x\in X$, and in
this case we say that the dynamical system $(\alg,\sfG,\alpha)$ is
\emph{$C_0(X)$-linear}.

\begin{Theorem}\label{thm:crossedfibres}
Let $X$ be a locally compact Hausdorff space, and let
$(\alg,\sfG,\alpha)$ be a $C_0(X)$-linear $C^*$-dynamical system. Then
the crossed product $\alg\rtimes_\alpha\sfG$ is again a
$C_0(X)$-algebra with fibres
\[
(\alg\rtimes_\alpha\sfG)_x\simeq\alg_x\rtimes_{\alpha^x}\sfG ,
\]
where
\[
\alpha_\gamma^x\big(a(x)\big) = \alpha_\gamma(a)(x)
\]
for each $x\in X$, $\gamma\in\sfG$ and $a\in\alg$.
\end{Theorem}

\begin{proof}The structure map
of $\alg\rtimes_\alpha\sfG$ is given by precomposing the structure map
of $\alg$ with the natural injection of the center of the multiplier
algebra of $\alg$ into the center of the multiplier algebra of
$\alg\rtimes_\alpha\sfG$; it satisfies
$(\mathtt{f}\cdot f)(\gamma)=\mathtt{f}\cdot\big(f(\gamma)\big)$ for all
$\mathtt{f}\in C_0(X)$, $f\in C_{\rm c}(\sfG,\alg)$ and
$\gamma\in\sfG$. See~\cite[Theorem~8.4]{Williams2007} for further
details.
\end{proof}


\setcounter{Theorem}{10}\setcounter{equation}{4}
% Original source lines 2534--2694
\begin{Remark}\label{rem:MoritaRKK}
There is a natural notion of Morita equivalence of $C^*$-algebra bundles over~$X$, similar to the notion of
equivariant Morita equivalence from Theorem~\ref{equivMorita}, which
uses the $C_0(X)$-bimodule structures: a~\emph{$C_0(X)$-linear Morita equivalence} between two
$C_0(X)$-algebras is a Morita equivalence which is compatible with the
$C^*$-bundle structures over~$X$. More generally, there is a category
$\mathcal{R}\CCK\hspace{-1mm}\CCK_X$ of $C^*$-algebra bundles over $X$ whose morphisms are
elements of Kasparov's groups $\mathcal{R}{\rm KK}(X;\alg,\balg)$, see, e.g., \cite{Echterhoff2009}: the cycles are the usual cycles
$(\hil,\phi,T)$ for Kasparov's bivariant K-theory ${\rm
 KK}(\alg,\balg)$ (cf.\ Section~\ref{sec:TdualityKK}) with the
additional requirement that $\phi\colon \alg\to\End_\balg(\hil)$ is
$C_0(X)$-linear. There is an obvious faithful functor
$\mathcal{R}\CCK\hspace{-1mm}\CCK_X \to \CCK\hspace{-1mm}\CCK$ which forgets the
$C_0(X)$-algebra structures. Isomorphic $C^*$-bundles in the category
$\mathcal{R}\CCK\hspace{-1mm}\CCK_X$ are precisely the $\mathcal{R}{\rm KK}$-equivalent
$C^*$-bundles. If $\alg$ and $\balg$ are isomorphic in
$\mathcal{R}\CCK\hspace{-1mm}\CCK_X$, i.e., there exists an invertible class
$\alpha\in \mathcal{R}{\rm KK}(X;\alg,\balg)$, then they are also
isomorphic as
$C^*$-algebras in the category $\CCK\hspace{-1mm}\CCK$.
\end{Remark}

\subsection[$\real^n$-actions on Mostow bundles]{$\boldsymbol{\real^n}$-actions on Mostow bundles}\label{sec:RMostow}

We can now apply the results of Section~\ref{sec:TopTtorus} to the
class of twisted tori given in
Section~\ref{sec:abeliansolv}.
The Mostow fibration of any almost abelian solvmanifold identifies
$\torus_{\Lambda_\sfG}$ as a torus bundle over a~circle, hence the
algebra of functions $C(\torus_{\Lambda_\sfG})$ is a~$C(\torus)$-algebra. In other words, $C(\torus_{\Lambda_\sfG})$ is an
object of the category $\mathcal{R}\CCK\hspace{-1mm}\CCK_\torus$, and we are
interested in the T-duality isomorphisms of~$C(\torus_{\Lambda_\sfG})$
in this category. In particular, given a fibrewise right action of the abelian
Lie group~$\R^n$ on~$\torus_{\Lambda_\sfG}$, it follows from
Theorem~\ref{thm:crossedfibres} that the $C^*$-algebraic T-dual
$C(\torus_{\Lambda_\sfG})\rtimes_\rt\R^n$ is also a $C(\torus)$-algebra, and by~\cite[Theorem~3.5]{Echterhoff2009} the $C^*$-bundles
$C(\torus_{\Lambda_\sfG})$ and
$C(\torus_{\Lambda_\sfG})\rtimes_\rt\R^n$ are isomorphic as $C^*$-algebras in the
category $\mathcal{R}\CCK\hspace{-1mm}\CCK_\torus$.

In order to have sensible definitions of
T-duality, we need to identify the
homologically non-trivial one-cycles of the twisted torus
$\torus_{\Lambda_\sfG}$, which are determined
in~\cite[Proposition~4.7]{Bock2009}. Write
\[
\ttM = (m_{ij})
\]
for the
integer matrix elements $m_{ij}\in\Z$ of the monodromy matrix. Since
$\sfG=\real^{d-1}\rtimes_\varphi\real$ is simply-connected, the
fundamental group of the twisted torus is
$\pi_1(\torus_{\Lambda_\sfG})\simeq\Lambda_\sfG$ whose abelianisation
$\Lambda_\sfG/[\Lambda_\sfG,\Lambda_\sfG]$ gives
the first homology group via the presentation
\begin{gather}\label{eq:H1torus}
{\rm H}_1(\torus_{\Lambda_\sfG},\zed) = \zed \oplus \bigg\langle
\hat e_1,\dots,\hat e_{d-1} \, \Big| \, \sum_{j=1}^{d-1} m_{ji} \hat
e_j = \hat e_i
\ \mbox{for} \
i=1,\dots,d-1 \bigg\rangle,
\end{gather}
where the first factor of $\zed$ corresponds to the base circle of the
torus fibration and the generators $\hat e_1,\dots,\hat e_{d-1}$
correspond to the torus fibres; they are given by
\begin{gather}\label{eq:hatei}
\hat e_i:=\sum_{k=1}^{d-1} \Sigma_{ki} \vec e_k,
\end{gather}
where
$\vec e_1,\ldots, \vec e_{d-1}$ is the standard basis of $\zed^{d-1}$ giving the
group law~\eqref{eq:grouplaw}.

We can then apply the structure theorem for finitely-generated
$\zed$-modules by appealing to some classical matrix algebra.
From the presentation \eqref{eq:H1torus}, ${\rm
 H}_1(\torus_{\Lambda_\sfG},\zed) = \zed \oplus {\rm{coker}}(\varphi_{x_0}-{\rm{id}}_{\Z^{d-1}})$
where $(\varphi_{x_0}-{\rm{id}}_{\Z^{d-1}})\colon \zed^{d-1}\to\zed^{d-1}$ in the
basis $\hat e_1,\ldots,\hat e_{d-1}$ is given by the integer relation matrix
\begin{gather}\label{eq:relationmatrix}
{\tt A}:=\ttM-\unit_{d-1}~.
\end{gather}
Let $r$ be the rank of $\tt A$. This matrix can be brought into its Smith normal form~$\tt D$
by finding invertible integer matrices ${\tt L}, {\tt
 R}\in\sfGL(d-1,\zed)$ such that
\[
{\tt D} = {\tt L} {\tt A} {\tt R}
\]
is diagonal with entries $m_i\in\zed$ for
$i=1,\dots,d-1$. The
integers~$m_i$ are the \emph{elementary divisors} of
${\tt A}$. They have the properties that $m_i$ divides
$m_{i+1}$, for $0<i<d-1$, and in particular $m_i=0$ for $i>r$; they can
be computed explicitly (up to sign) as
\[
m_i = \frac{d_i({\tt A})}{d_{i-1}({\tt A})},
\]
where the \emph{$i$-th determinant divisor} $d_i({\tt A})$ is the greatest common
divisor of all $i{\times} i$ minors of the relation matrix ${\tt A}$,
with $d_0({\tt A}):=1$. The matrices ${\tt L}, {\tt
 R}\in\sfGL(d-1,\zed)$ are found by reducing the matrix ${\tt A}$ to
its Smith normal form ${\tt D}$ through a sequence of elementary row
and column operations over~$\zed$, see, e.g.,~\cite{Havas1997}.

Given the Smith normal form, we set
\begin{gather}\label{eq:tildeei}
\tilde e_i:=\sum_{j=1}^{d-1}\,\big({\tt L}^{-1}\big)_{ji}\,\hat e_j
\end{gather}
and observe that
the image of $\tt A$, which is generated over $\zed$ by the vectors
\[
\sum_{j=1}^{d-1} {\tt A}_{ji} \hat e_j=\sum_{k,l=1}^{d-1} \big({\tt
 R}^{-1}\big)_{ki} {\tt
 D}_{lk} \tilde e_k ,
\]
is equivalently generated by the
vectors $m_k \tilde e_k$ with $k=1,\ldots, r$.
Hence
\begin{align*}
{\rm{coker}}(\varphi_{x_0}-{\rm{id}}_{\Z^{d-1}})&=\big\langle \hat e_1, \ldots , \hat
e_{d-1}\big\rangle\Big/\bigg\langle \sum\limits_{j=1}^{d-1} {\tt A}_{j1} \hat e_j,\ldots,
 \sum\limits_{j=1}^{d-1} {\tt A}_{j\:\!d-1} \hat e_j\bigg\rangle\\
&=\big\langle \tilde e_1,
\ldots , \tilde e_{d-1}\big\rangle\big/\big\langle m_1 \tilde e_1,\ldots, m_r \tilde
e_r\big\rangle
\end{align*}
and
\begin{gather}\label{eq:H1gendecomp}
{\rm H}_1(\torus_{\Lambda_\sfG},\zed) \simeq \zed \oplus
\zed^{d-1-r} \oplus \bigoplus_{i=1}^r \zed_{m_i} .
\end{gather}

We are exclusively
interested in the natural $\R^n$-actions on $\torus_{\Lambda_\sfG}$ which descend
from actions of abelian subgroups $\R^n\subset\sfG$, acting on $\sfG$
by right multiplication. They can be
organised into three classes associated with the different types of
summands in the $\Z$-module presentation of the homology group~\eqref{eq:H1gendecomp}, and we only retain those which are fiberwise
actions on the Mostow bundle.
The first summand~$\Z$ in~\eqref{eq:H1gendecomp} corresponds to the subgroup
\[
\R_x = \big\{(0,\xi)\in\sfG\big\}
\]
acting on $\sfG$ by right multiplication:
\[
(z,x) (0,\xi) = (z,x+\xi)
\]
for all $(z,x)\in\sfG$ and $\xi\in\R$. Clearly this does not descend
to a fiberwise
action on the twisted torus $\torus_{\Lambda_\sfG}$, and the crossed
product $C(\torus_{\Lambda_\sfG})\rtimes_\rt\R_x$ is no longer a
$C(\torus)$-algebra. Thus an $\R_x$-action takes us out of the
category $\mathcal{R}\CCK\hspace{-1mm}\CCK_\torus$, and we will henceforth
discard $\R^n$-actions where $\R^n$ contains the subgroup~$\R_x$.

For the remaining types of summands in \eqref{eq:H1gendecomp}, we can
give explicit descriptions of the $C^*$-algebraic T-duals of an almost
abelian solvmanifold. We consider both classes in turn. As the only solvmanifolds in one and two dimensions are tori, which are
already treated by our analysis from
Section~\ref{sec:toruscrossedprod}, we assume $d\geq3$ for the
remainder of this paper.


\setcounter{subsection}{5}\setcounter{Theorem}{15}\setcounter{equation}{9}
% Original source lines 2870--3217
\subsection[$\real_z$-actions: noncommutative torus bundles]{$\boldsymbol{\real_z}$-actions: noncommutative torus bundles}\label{sec:Rzgen}

Let us now come to the torsion summands $\Z_{m_i}$ in the
$\Z$-module decomposition \eqref{eq:H1gendecomp}. Pick a~non-trivial
elementary divisor $m_i>1$ for some $i\in\{1,\dots,r\}$. The
corresponding homology generator $\tilde e_i$ is constructed as a~$\Z$-linear combination~\eqref{eq:tildeei} of the generators~\eqref{eq:hatei}. It defines a~fixed
element~$z_0$ in the image of $\varphi_{x_0}-\rm{id}_{\R^{d-1}}$ and a~corresponding one-dimensional subgroup of~$\sfG$ given by
\[
\R_{z_0} := \R (z_0,0) .
\]
We further assume that
\[
\Z_{ z_0} := \R_{ z_0}\cap\Lambda_\sfG
\]
is a lattice in $\R_{z_0}$.
The choice of $z_0\in\R^{d-1}$ is not unique, and
any change of basis of $\Z^{d-1}$, represented by a matrix
$B\in\sfGL(d-1,\Z)$, defines an equally good element $B\cdot
z_0\in\R^{d-1}$ as long as $B\cdot
z_0\in{\rm im}\big(\varphi_{x_0}-\Id_{\R^{d-1}}\big)$. We write $\proj_{z_0}$ for the
linear projection of $\R^{d-1}$ to $\R_{z_0}$, and denote by $\langle
z_0,z\rangle\in\R$ the component of $z\in\R^{d-1}$ in $\R_{z_0}$,
i.e.,~$\proj_{z_0}\cdot z = \langle z_0,z\rangle \, z_0$.

In the basis $\hat e_i$, the lattice of $\sfG$ is given by
\[
\Lambda_\sfG=\Z^{d-1}\rtimes_{\hat\varphi|_{x_0\,\Z}} x_0 \Z ,
\]
where $\hat\varphi_x:=\Sigma^{-1}\,\varphi_x\,\Sigma$ for all
$x\in\R$. Then the fibres of the underlying Mostow bundle are `square' tori
$\torus^{d-1}$ with unit periodicities $ z\sim z+\vec
e_i$ for $i=1,\dots,d-1$, where as
before $\vec e_i$ denotes the standard basis of
$\R^{d-1}$. The action of the
subgroup $\R_{ z_0}$ on
elements $(z,x)\in\sfG$ by right multiplication is given by
\begin{gather}\label{eq:Rz0action}
(z,x) (\zeta\,z_0,0) = \big(z+\zeta \Sigma^{-1} \varphi_x\cdot
z_0,x\big)
\end{gather}
for $\zeta\in\R$, where $\Sigma^{-1} \varphi_x\cdot z_0$ lies in the
image of the relation matrix ${\tt A}=\ttM-\unit_{d-1}$.

Our principal tool to compute the $C^*$-algebraic
T-dual for such an action of $\R$ in the image of $\varphi_{x_0}-\Id_{\R^{d-1}}$, which we collectively refer to as
$\R_z$-actions, will be Green's theorem in the form~\eqref{scheme}:
\begin{gather}\label{Mostowscheme}
C(\torus_{\Lambda_\sfG})\rtimes_\rt \real_{ z_0} \ \mbf\sim_{\text{\tiny M}} \
C_0(\sfG/\real_{ z_0})\rtimes_\lt \Lambda_\sfG .
\end{gather}
By appealing to
Theorem~\ref{thm:crossedfibres}, we may apply \eqref{Mostowscheme}
fibrewise. For fixed $x\in\R$, the fibre $\sfG_x$ of the semi-direct
product $\sfG=\R^{d-1}\rtimes_{\hat\varphi}\R$ is the subgroup~$\R^{d-1}$, and the corresponding fibre of
the solvmanifold $\torus_{\Lambda_\sfG}$ over $x\in\R/x_0\,\Z$ is the
torus $\torus^{d-1}=\R^{d-1}/\Z^{d-1}$. The Morita
equivalence~\eqref{Mostowscheme} is $C(\torus)$-linear and the fibres of the corresponding
T-dual $C(\torus)$-algebra are given by the fibrewise Morita
equivalence
\begin{gather}\label{eq:Mostowschemefiber}
\big(C(\torus_{\Lambda_\sfG})\rtimes_\rt\real_{ z_0}\big)_x \simeq
C\big(\R^{d-1}/\Z^{d-1}\big)\rtimes_{\rt^x}\real_{ z_0} \
\mbf\sim_{\text{\tiny M}} \
C_0\big(\real^{d-1}/\R_{ z_0}\big)\rtimes_{\lt^x}\Z^{d-1} .
\end{gather}
The action of the subgroup $\Z^{d-1}\subset\Lambda_\sfG$ on the
coset space $\R^{d-1}/\R_{ z_0}$ is induced by left multiplication in the
group $\sfG$.
After a basis transformation, we can decompose the discrete group $\Z^{d-2}$ into a direct sum
\[
\Z^{d-1} \simeq \Z_v^{d-2}\oplus\Z_{ z_0} ,
\]
where $\Z_v^{d-1} = (\unit_{d-1}-\proj_{z_0})\cdot\Z^{d-1}$.

Let
\[
\mathbb{F}_* := \big\{x\in\R \, \big| \, \big\langle
z_0,\Sigma^{-1} \varphi_x\cdot z_0\big\rangle=0\big\} .
\]
\begin{Lemma}\label{lem:ellfibercom}
Over any $x\in \mathbb{F}_*$, the fiber
$\big(C(\torus_{\Lambda_\sfG})\rtimes_\rt\real_{ z_0}\big)_x$
is Morita equivalent to the
commutative $C^*$-algebra \smash{$C\big(\torus^{d-1}\big)$}.
\end{Lemma}
\begin{proof}
If $x\in\mathbb{F}_*$, then $w_0:=\Sigma^{-1}\,\varphi_x\cdot z_0$ only
shifts the corresponding component of $ z$ in
$(\unit_{d-1}-\proj_{ z_0})\cdot\R^{d-1}$ in the $\R_{z}$-action
\eqref{eq:Rz0action}. In this case any
element $(z,x)\in\sfG$ may be factorized as
$$
(z,x) = \big(\proj_{ z_0}\cdot
z+(\unit_{d-1}-\proj_{z_0}-\proj_{w_0})\cdot z\,,\,x\big) \,
\big(\langle w_0,z\rangle\, z_0 \,,\,0\big) \ .
$$
Hence the coset space $\R^{d-1}/\R_{ z_0}\simeq\R^{d-2}$ can be
parameterized by $\proj_{ z_0}\cdot
z+\big(\unit_{d-1}-\proj_{z_0}-\proj_{w_0}\big)\cdot z$ with $z\in\R^{d-1}$,
which we decompose correspondingly into a direct product
$\R_{z_0}\times\R^{d-3}$ (with the second factor absent for $d=3$). By \eqref{eq:Mostowtorus} the discrete group
$\Z_v^{d-2}$ acts trivially on the line $\R_{z_0}$ and by translations
on $\R^{d-3}$,
while~$\Z_{ z}$ acts by translations on $\R_{z_0}$ and trivially
on $\R^{d-3}$. Then the crossed product on the right-hand side of~\eqref{eq:Mostowschemefiber} may be unravelled to get
\begin{align*}
C_0\big(\real^{d-1}/\R_{ z_0}\big)\rtimes_{\lt^x}\Z^{d-1}
&\simeq C_0\big(\R_{
 z_0}\times\R^{d-3}\big)\rtimes_{\lt^x}\big(\Z_v^{d-2}\times\Z_{ z_0}\big) \\
&\simeq
\big[\big(C_0(\R_{z_0})\otimes C_0\big(\R^{d-3}\big)\big)\rtimes_{\Id\otimes\lt^x}\Z^{d-2}_v\big]\rtimes_{\beta^x}\Z_{z_0} \\
& \ \mbf\sim_{\text{\tiny M}} \ \big(C_0(\R_{ z_0})\otimes C^*(\Z)\otimes C\big(\torus^{d-3}\big)\big)\rtimes_{\lt^x\otimes\Id\otimes\Id}\Z_{ z_0}
 \\
&\simeq \big(C_0(\R_{ z_0})\rtimes_{\lt^x}\Z_{ z_0}\big)\otimes
C^*(\Z)\otimes C\big(\torus^{d-3}\big)
\nonumber\\
& \ \mbf\sim_{\text{\tiny M}} \ C(\torus)\otimes C^*(\Z)\otimes C\big(\torus^{d-3}\big) \simeq C\big(\torus^{d-1}\big) .
\end{align*}
In the second line we applied
Theorem~\ref{cpsp3}. In the third line we used
Example~\ref{ex:Moritatori} together with the fact that the
homomorphism $\sigma_{\Z_{ z_0}}$ from \eqref{sigmaI} is trivial since the
groups $\Z_{ z_0}$ and $\Z_v^{d-2}$ are discrete, and so the action of $\Z_{ z_0}$ on the crossed
product is induced by left multiplication on $\R_{ z_0}$, the trivial
action on $\R^{d-3}$, and the
trivial action on $\Z_v^{d-2}$. In the fifth line we used
Example~\ref{ex:Moritatori} again.
\end{proof}

The central result of this paper is
\begin{Theorem}
\label{thm:NCtorusbundlegen}
The $C^*$-algebraic T-dual of any almost abelian
solvmanifold $\torus_{\Lambda_\sfG}$ with respect to an $\R_z$-action
is Morita equivalent to a $C^*$-algebra bundle
of noncommutative tori over the circle $\torus$:
\begin{gather*}%\label{eq:algbundlegen}
C(\torus_{\Lambda_\sfG})\rtimes_\rt\real_{ z_0}
 \ \mbf\sim_{\text{\tiny M}} \ \mcoprod
\torus^{d-1}_{\vec\theta_{z_0}(x)} ,
\end{gather*}
where the noncommutativity parameters $\vec\theta_{z_0}(x)\in\R^{d-2}$ are given
by
\begin{gather}\label{eq:thetaz0}
\vec\theta_{z_0}(x) = \begin{cases}
 0 & \mbox{for} \ x \in \F_* , \\
 \dfrac{(\unit_{d-1}-\proj_{ z_0}) \Sigma^{-1} \varphi_{x_0 x}\cdot
 z_0}{\big\langle{ z_0},\Sigma^{-1} \varphi_{x_0 x}\cdot z_0\big\rangle}
 & \mbox{for} \ x\in\R\setminus\F_* . \end{cases}
\end{gather}
\end{Theorem}
\begin{proof}
That the fibres over $x\in\F_*$ are just ordinary tori $\torus^{d-1}$
is established by Lemma~\ref{lem:ellfibercom}, so we may assume that $x\in\R\setminus\F_*$. Then a simple
calculation shows that any element $(z,x)\in\sfG$ can be factorized as
\[
(z,x) = (v,x) \left(\frac{\proj_{ z_0}\cdot z}{\big\langle z_0,\Sigma^{-1} \varphi_x\cdot z_0\big\rangle} , 0\right) ,
\]
where
\begin{gather}\label{eq:Mostowcoset}
v:=\big(\unit_{d-1}-\proj_{ z_0}\big)\cdot z - \frac{\langle{ z_0},
 z\rangle}{\big\langle{ z_0},\Sigma^{-1} \varphi_x\cdot z_0\big\rangle}
\big(\unit_{d-1}-\proj_{ z_0}\big) \Sigma^{-1} \varphi_x\cdot z_0 .
\end{gather}
Thus the coset space $\R_{x,v}^{d-2}:=\R^{d-1}/\R_{ z_0}$ may be
parameterized by the coordinates $v\in\R^{d-2}$ over any
$x\in\R\setminus\F_*$, and we explicitly retain the fibre index in the
notation for convenience.

We now need to unravel the crossed product
on the right-hand side of~\eqref{eq:Mostowschemefiber}.
From~\eqref{eq:Mostowtorus} and~\eqref{eq:Mostowcoset} it follows that the action of elements
$(\gamma_v,\gamma_{ z_0})\in\Z^{d-1}\simeq\Z_v^{d-2}\oplus\Z_{
 z_0}$ on the coset is given by
\begin{gather}\label{eq:Zcosetactionv}
(\gamma_v,0,0)\cdot (v,x) = (v+\gamma_v,x) , \\
(0,\gamma_{ z_0},0)\cdot (v,x) = \left(v-\frac{
 \gamma_{ z_0}}{\big\langle{ z_0},\Sigma^{-1} \varphi_x\cdot z_0\big\rangle}
\big(\unit_{d-1}-\proj_{ z_0}\big) \Sigma^{-1} \varphi_x\cdot
 z_0 , x\right) .
\label{eq:Zcosetactionz}\end{gather}
Applying
Theorem~\ref{cpsp3} as in the proof of Lemma~\ref{lem:ellfibercom} gives
\begin{gather*}%\label{eq:MostowRHS}
C\big(\R^{d-1}/\Z^{d-1}\big)\rtimes_{\rt^x}\real_{ z_0}
\ \mbf\sim_{\text{\tiny M}} \
\big(C_0\big(\real^{d-2}_{x,v}\big)\rtimes_{\lt^x}\Z^{d-2}_{v}\big) \rtimes_{\beta^x} \Z_{ z_0} ,
\end{gather*}
where the action of $\Z^{d-2}_v$ on the coset $\R^{d-2}_{x,v}$ is
given by~\eqref{eq:Zcosetactionv}, while the action of $\Z_{ z_0}$ on the crossed product
$C_0\big(\real^{d-2}_{x,v}\big)\rtimes_{\lt^x}\Z^{d-2}_{v}$ is induced by the
action on $\real^{d-2}_{x,v}$ given in~\eqref{eq:Zcosetactionz} and the trivial action on~$\Z^{d-2}_{v}$.

Next, Example~\ref{ex:Moritatori} yields
\begin{gather}\label{eq:NMoritav}
C_0\big(\real_{x,v}^{d-2}\big)\rtimes_{\lt^x} \Z_{v}^{d-2} \ \mbf\sim_{\text{\tiny M}} \
C\big(\torus^{d-2}\big) .
\end{gather}
Since the homomorphism
$\sigma_{\Z_{ z_0}}\colon \Z_{ z_0}\to\real^+$ is trivial, it is not difficult to see
that the Morita equivalence bimodule $\Mcal$ implementing this
equivalence is $\Z_{ z_0}$-equivariant (in the sense of Theorem~\ref{equivMorita}):
the $\Z_{ z_0}$-action
$U\colon \Z_{ z_0}\rightarrow \Aut(\Mcal)$ is given by $U_{\gamma_v}(\xi)(v)=\xi(v+\gamma_v)$, for all $\gamma_v\in
\Z_{ z_0}$, $\xi\in \Mcal$ and $v\in \real_{x,v}^{d-2} $.
Applying Theorem \ref{equivMorita} we conclude that the
Morita equivalence \eqref{eq:NMoritav} induces the Morita equivalence
\begin{gather}\label{eq:fibreNCMostow}
\big(C_0\big(\real^{d-2}_{x,v}\big)\rtimes_{\lt^x}\Z^{d-2}_{v}\big) \rtimes_{\beta^x}
\Z_{ z_0} \ \mbf\sim_{\text{\tiny M}} \
C\big(\torus^{d-2}\big)\rtimes_{\lt^x}\Z_{ z_0} ,
\end{gather}
where the action
of the group $\Z_{ z_0}$ on $C\big(\torus^{d-2}\big)$ is the pullback of the action on
$\torus^{d-2}=\real_{x,v}^{d-2}/\Z_{v}^{d-2}$ induced
from~\eqref{eq:Zcosetactionz}.

After rescaling $x$ by $x_0$
to give it unit period, this shows that the
corresponding algebra of functions on the fiber at $\e^{2\pi\ii x}$ is that
of a noncommutative $d{-}1$-torus $\torus_{\vec\theta_{
 z_0}(x)}^{d-1}$ with noncommutativity parameter
$\vec\theta_{z_0}(x)$ given by~\eqref{eq:thetaz0}, see Example~\ref{ex:NCtorid}.
Thus the $C^*$-algebraic T-dual
$C(\torus_{\Lambda_\sfG})\rtimes_\rt\real_{ z_0}$ is a $C^*$-algebra bundle
of noncommutative tori $\A_{\vec\theta_{
 z_0}(x)}=\torus^{d-1}_{\vec\theta_{ z_0}(x)}$ over the
circle $\torus=\big\{\e^{2\pi\ii x}\, | \, x\in \real/\Z\big\}$.
\end{proof}

What becomes of the non-trivial monodromy \eqref{eq:Mostowtorus} of
the $\torus^{d-1}$ fibers parameterized by~$z$ in the original Mostow
bundle? To answer this question, we write the monodromy matrix
$\ttM=(m_{ij})\in\sfSL(d-1,\Z)$ in the block form
\begin{gather}\label{eq:ttMblock}
\ttM = \left(\begin{matrix} \ttM_{|d-2} & \vec m \\
\vec m' & m_{d-1 \, d-1}
\end{matrix}\right) ,
\end{gather}
where $\ttM_{|d-2}=(m_{ij})_{1\leq i,j\leq d-2}$, while $\vec m=(m_{i \,
 d-1})_{i=1,\dots,d-2}$ and $\vec m' = (m_{d-1 \,
 i})_{i=1,\dots,d-2}$ are respectively column and row vectors in $\Z^{d-2}$. Below
we denote the usual Euclidean inner product on $\R^{d-2}$ by
$\langle\,\cdot\,,\,\cdot\,\rangle$.
\begin{Proposition}\label{prop:monodromyprop}
The noncommutativity parameter
$\vec\theta_{\vec e_{d-1}}(x)$ from~\eqref{eq:thetaz0} varies with a
change of coset representative $x\in\R/\Z$ according to its
$\sfSL(d-1,\Z)$ orbit under the action of the monodromy matrix
\eqref{eq:ttMblock} by $(d{-}2)$-dimensional linear fractional transformations
\begin{gather}\label{eq:thetaMoritagen}
\vec\theta_{\vec e_{d-1}}(x+1) = \ttM\big[\vec\theta_{\vec
 e_{d-1}}(x)\big] :=
\frac{\ttM_{|d-2}\cdot\vec\theta_{\vec e_{d-1}}(x) + \vec m}
{\big\langle \vec m',\vec\theta_{\vec e_{d-1}}(x)\big\rangle + m_{d-1 \,
 d-1}} .
\end{gather}
Under a change of basis of $\Z^{d-1}$ given
by a matrix $B\in\sfGL(d-1,\Z)$, the corresponding noncommutativity
parameter varies according to
\[
\vec\theta_{B\cdot\vec e_{d-1}}(x+1) = \big(B^{\rm
 t} \ttM B\big)\big[\vec\theta_{B\cdot\vec e_{d-1}}(x)\big] .
\]
\end{Proposition}
\begin{proof}
The key feature stems from the definition \eqref{eq:monodromydef} of the monodromy matrix
$\ttM$:
\[
\Sigma^{-1}\,\varphi_{x+x_0} = \Sigma^{-1} \varphi_{x_0} \varphi_x =
\ttM \Sigma^{-1} \varphi_x .
\]
The statements then follow from straightforward calculations in components.
\end{proof}

\begin{Remark}
\label{rem:MoritaKK}
The right-hand side of \eqref{eq:thetaMoritagen} is an example of a
linear fractional transformation in higher dimensions, known from
complex analysis, see, e.g.,~\cite{Cowen2000}. We can show that it defines a Morita equivalence of noncommutative $d{-}1$-tori from
Example~\ref{ex:NCTdMorita}. For this, we introduce the skew-symmetric
matrix $\Theta$ corresponding to the vector $\vec\theta\in\R^{d-2}$ as
in Example~\ref{ex:NCtorid}:
\[
\Theta = \left(
\begin{matrix}
\mbf 0_{d-2} & \vec\theta \\ -\vec\theta\,^{\rm t} & 0
\end{matrix}
\right) .
\]
We denote by $\mbf g_{\ttM}$ the element of ${\sf SO}(d-1,d-1;\Z)$
corresponding to the monodromy matrix $\ttM\in\sfSL(d-1,\Z)$:
\[
\mbf g_{\ttM} = \left(
\begin{matrix}
\ttM & \mbf 0_{d-1} \\ \mbf 0_{d-1} & \big(\ttM^{\rm t}\big)^{-1}
\end{matrix}
\right)
\in {\sf SO}(d-1,d-1;\Z) .
\]
Introduce matrices $\mbf T_{d-1}$ of order $2$ with determinant $-1$
by
\[
\mbf T_{d-1} = \left(
\begin{matrix}
\unit_{d-1}-{\tt E}_{d-1} & {\tt E}_{d-1} \\ {\tt E}_{d-1} &
\unit_{d-1}-{\tt E}_{d-1}
\end{matrix}
\right) \in {\sf O}(d-1,d-1;\Z) ,
\]
where ${\tt E}_{d-1}$ is the matrix unit whose only
non-zero element is $({\tt E}_{d-1})_{d-1\,d-1} = 1$.
Finally, define the element $\mbf M\in{\sf SO}(d-1,d-1;\Z)$ by
\[
\mbf M = \mbf T_{d-1} \mbf g_{\ttM} \mbf T_{d-1} .
\]
Using the adjugate formula
for matrix inverses, a straightforward if tedious calculation then shows that the
corresponding ${\sf SO}(d-1,d-1;\Z)$ orbit of $\Theta$ reproduces the
$(d{-}2)$-dimensional linear fractional transformation of
\eqref{eq:thetaMoritagen}:
\[
\mbf M[\Theta] = \left(
\begin{matrix}
\mbf 0_{d-2} & \ttM\big[\vec\theta\, \big] \\
-\ttM\big[\vec\theta\, \big]^{\rm t} & 0
\end{matrix}
\right) \qquad \mbox{with} \quad \ttM\big[\vec\theta\, \big] = \frac{\ttM_{|d-2}\cdot\vec\theta + \vec m}
{\big\langle \vec m',\vec\theta\, \big\rangle + m_{d-1 \,
 d-1}} .
\]
It follows that, while the fibre noncommutative
tori of Theorem~\ref{thm:NCtorusbundlegen} are not generally identical under a change of
representative $x\in\R/\Z$, by Example~\ref{ex:NCTdMorita}
they are always Morita
equivalent:
\[
\torus^{d-1}_{\vec\theta_{\vec e_{d-1}}(x+1)} \ \mbf\sim_{\text{\tiny
 M}} \ \torus^{d-1}_{\vec\theta_{\vec e_{d-1}(x)}} .
\]
Recalling that our formulation of topological T-duality takes place in the additive category~${\CCK\hspace{-1mm}\CCK}$ from
Section~\ref{sec:TdualityKK}, this has a natural interpretation: The
non-trivial monodromy in the automorphism group of the~\smash{$\torus^{d-1}$} fibres of the
twisted torus $\torus_{\Lambda_\sfG}$ is
manifested as a (generally non-trivial) isomorphism in the Morita automorphism
group of the fibre noncommutative tori in~${\CCK\hspace{-1mm}\CCK}$.
\end{Remark}

\begin{thebibliography}{99}
\bibitem[2]{Bock2009}
Bock C., On low-dimensional solvmanifolds, \href{https://doi.org/10.4310/AJM.2016.v20.n2.a1}{\textit{Asian~J. Math.}} \textbf{20}
 (2016), 199--262, \href{https://arxiv.org/abs/0903.2926}{arXiv:0903.2926}.

\bibitem[6]{Brodzki2006}
Brodzki J., Mathai V., Rosenberg J., Szabo R.J., D-branes, {RR}-fields and
 duality on noncommutative manifolds, \href{https://doi.org/10.1007/s00220-007-0396-y}{\textit{Comm. Math. Phys.}} \textbf{277}
 (2008), 643--706, \href{https://arxiv.org/abs/hep-th/0607020}{arXiv:hep-th/0607020}.

\bibitem[7]{Brodzki2007}
Brodzki J., Mathai V., Rosenberg J., Szabo R.J., Non-commutative
 correspondences, duality and {D}-branes in bivariant {$K$}-theory,
 \href{https://doi.org/10.4310/ATMP.2009.v13.n2.a4}{\textit{Adv. Theor. Math. Phys.}} \textbf{13} (2009), 497--552,
 \href{https://arxiv.org/abs/0708.2648}{arXiv:0708.2648}.

\bibitem[10]{Console2016}
Console S., Macr\`{\i} M., Lattices, cohomology and models of 6-dimensional
 almost abelian solvmanifolds, \textit{Rend. Semin. Mat. Univ. Politec.
 Torino} \textbf{74} (2016), 95--119, \href{https://arxiv.org/abs/1206.5977}{arXiv:1206.5977}.

\bibitem[11]{Corwin1990}
Corwin L.J., Greenleaf F.P., Representations of nilpotent {L}ie groups and
 their applications. {P}art~{I}. Basic theory and examples, \textit{Cambridge
 Studies in Advanced Mathematics}, Vol.~18, Cambridge University Press,
 Cambridge, 1990.

\bibitem[12]{Cowen2000}
Cowen C.C., MacCluer B.D., Linear fractional maps of the ball and their
 composition operators, \textit{Acta Sci. Math. (Szeged)} \textbf{66} (2000),
 351--376.

\bibitem[14]{Echterhoff2017}
Echterhoff S., Crossed products and the Mackey--Rieffel--Green machine,
 \href{https://doi.org/10.1088/1126-6708/2003/09/054}{\textit{Oberwolfach Sem.}} \textbf{47} (2017), 5--79, \href{https://arxiv.org/abs/1006.4975}{arXiv:1006.4975}.

\bibitem[15]{Echterhoff2009}
Echterhoff S., Nest R., Oyono-Oyono H., Principal non-commutative torus
 bundles, \href{https://doi.org/10.1112/plms/pdn050}{\textit{Proc. Lond. Math. Soc.}} \textbf{99} (2009), 1--31,
 \href{https://arxiv.org/abs/#2}{arXiv:0810.0811}.

\bibitem[17]{Fack1981}
Fack T., Skandalis G., Connes' analogue of the {T}hom isomorphism for the
 {K}asparov groups, \href{https://doi.org/10.1007/BF01393931}{\textit{Invent. Math.}} \textbf{64} (1981), 7--14.

\bibitem[23]{Havas1997}
Havas G., Majewski B.S., Integer matrix diagonalization, \href{https://doi.org/10.1006/jsco.1996.0141}{\textit{J.~Symbolic
 Comput.}} \textbf{24} (1997), 399--408.

\bibitem[32]{Mathai2004}
Mathai V., Rosenberg J., {$T$}-duality for torus bundles with {$H$}-fluxes via
 noncommutative topology, \href{https://doi.org/10.1007/s00220-004-1159-7}{\textit{Comm. Math. Phys.}} \textbf{253} (2005),
 705--721, \href{https://arxiv.org/abs/hep-th/0401168}{arXiv:hep-th/0401168}.

\bibitem[34]{Milnor1976}
Milnor J., Curvatures of left invariant metrics on {L}ie groups, \href{https://doi.org/10.1016/S0001-8708(76)80002-3}{\textit{Adv.
 Math.}} \textbf{21} (1976), 293--329.

\bibitem[35]{Mostow1954}
Mostow G.D., Factor spaces of solvable groups, \href{https://doi.org/10.2307/1969700}{\textit{Ann. of Math.}}
 \textbf{60} (1954), 1--27.

\bibitem[37]{Phillips2006}
Phillips N.C., Every simple higher-dimensional noncommutative torus is an AT
 algebra, \href{https://arxiv.org/abs/math.OA/0609783}{arXiv:math.OA/0609783}.

\bibitem[40]{Raghunathan1972}
Raghunathan M.S., Discrete subgroups of {L}ie groups, \textit{Ergebnisse der
 Mathematik und ihrer Grenzgebiete}, Vol.~68, Springer-Verlag, New York~--
 Heidelberg, 1972.

\bibitem[41]{Rieffel1974}
Rieffel M.A., Induced representations of {$C^{\ast}$}-algebras, \href{https://doi.org/10.1016/0001-8708(74)90068-1}{\textit{Adv.
 Math.}} \textbf{13} (1974), 176--257.

\bibitem[42]{Rieffel1976}
Rieffel M.A., Strong {M}orita equivalence of certain transformation group
 {$C^*$}-algebras, \href{https://doi.org/10.1007/BF01418238}{\textit{Math. Ann.}} \textbf{222} (1976), 7--22.

\bibitem[43]{Rieffel1981}
Rieffel M.A., {$C^{\ast} $}-algebras associated with irrational rotations,
 \href{https://doi.org/10.2140/pjm.1981.93.415}{\textit{Pacific~J. Math.}} \textbf{93} (1981), 415--429.

\bibitem[44]{Rieffel1982}
Rieffel M.A., The cancellation theorem for projective modules over irrational
 rotation {$C^{\ast} $}-algebras, \href{https://doi.org/10.1112/plms/s3-47.2.285}{\textit{Proc. London Math. Soc.}} \textbf{47}
 (1983), 285--302.

\bibitem[45]{Rieffel1990}
Rieffel M.A., Noncommutative tori~-- a case study of noncommutative
 differentiable manifolds, in Geometric and Topological Invariants of Elliptic
 Operators ({B}runswick, {ME}, 1988), \textit{Contemp. Math.}, Vol.~105, \href{https://doi.org/10.1090/conm/105/1047281}{Amer.
 Math. Soc.}, Providence, RI, 1990, 191--211.

\bibitem[46]{Rieffel1998}
Rieffel M.A., Schwarz A., Morita equivalence of multidimensional noncommutative
 tori, \href{https://doi.org/10.1142/S0129167X99000100}{\textit{Internat.~J. Math.}} \textbf{10} (1999), 289--299,
 \href{https://arxiv.org/abs/math.QA/9803057}{arXiv:math.QA/9803057}.

\bibitem[50]{OpreaTralle}
Tralle A., Oprea J., Symplectic manifolds with no {K}\"ahler structure,
 \textit{Lecture Notes in Mathematics}, Vol.~1661, \href{https://doi.org/10.1007/BFb0092608}{Springer-Verlag}, Berlin,
 1997.

\bibitem[52]{Williams2007}
Williams D.P., Crossed products of {$C{^\ast}$}-algebras, \textit{Mathematical
 Surveys and Monographs}, Vol.~134, \href{https://doi.org/10.1090/surv/134}{Amer. Math. Soc.}, Providence, RI, 2007.

\end{thebibliography}
\end{document}
